% 第六讲：从需求数据到偏好
% 第 6 章 统一书稿前言包装
\newcommand{\lectureno}{第六章}
\newcommand{\lectureshort}{从需求数据到偏好}
% ============================================================
%  微观经济学理论课程讲义 —— 统一课程前言（四讲共用）
%  版式参考 Nolan Miller, Notes on Microeconomic Theory (SWP/notes3 风格)
%  编译：LuaLaTeX（两遍）
%  各讲 preamble.tex 只需设定 \lectureno / \lectureshort 后 \input 本文件。
% ============================================================
\documentclass[11pt,a4paper,fontset=none,oneside,openany]{ctexbook}

% ---------- 版式 ----------
\usepackage{geometry}
\geometry{
  a4paper,
  top=2.8cm, bottom=2.8cm,
  left=2.6cm, right=2.6cm,
  headsep=0.85cm, footskip=1.25cm
}
\setlength{\headheight}{14pt}

% ---------- 字体（全部矢量嵌入，保证 PDF 高清晰度）----------
% 注意：本机 macOS 已移除 Kaiti SC，CJK 斜体统一回退到 Songti SC；
%        Songti SC / Heiti SC 经 OSFONTDIR 指向 Supplemental 后由 luaotfload 解析。
\usepackage{fontspec}
\setmainfont{TeX Gyre Termes}[Ligatures=TeX]
\setsansfont{TeX Gyre Heros}[Ligatures=TeX,Scale=MatchLowercase]
\setmonofont{TeX Gyre Cursor}[Scale=MatchLowercase]
\setCJKmainfont{Songti SC}[BoldFont={Songti SC Bold},ItalicFont={Songti SC}]
\setCJKsansfont{Heiti SC}[BoldFont={Heiti SC}]

% ---------- 数学 ----------
\usepackage{amsmath}
\usepackage{mathtools}
\usepackage{amsthm}
\usepackage{unicode-math}
\setmathfont{TeX Gyre Termes Math}
\allowdisplaybreaks[2]
\setlength{\abovedisplayskip}{9pt plus 3pt minus 3pt}
\setlength{\belowdisplayskip}{9pt plus 3pt minus 3pt}

% ---------- 定理类环境（书稿格式：按章编号，如“定理 1.1”）----------
\theoremstyle{plain}
\newtheorem{theorem}{定理}[chapter]
\newtheorem{proposition}[theorem]{命题}
\newtheorem{lemma}[theorem]{引理}
\newtheorem{corollary}[theorem]{推论}
\theoremstyle{definition}
\newtheorem{definition}[theorem]{定义}
\newtheorem{condition}[theorem]{条件}
\newtheorem{example}[theorem]{例}
\theoremstyle{remark}
\newtheorem{remark}[theorem]{注}
% 公式、图、表按章编号（书稿格式：式 (1.1)、图 1.1、表 1.1）
\numberwithin{equation}{chapter}
\numberwithin{figure}{chapter}
\numberwithin{table}{chapter}

% 定理类环境标题后的标点：amsthm 默认加英文 "."，此处统一改为中文句号。
% 原理：amsthm 的 \@thm 每次开环境都执行 \thm@headpunct{.}，而样式宏在其之后展开，
% 故把 \thm@headpunct{。} 追加到三个样式宏末尾即可覆盖（用 \AtBeginDocument 以免被 ctex 覆盖）。
\makeatletter
\AtBeginDocument{%
  \let\WBth@plain\th@plain
  \let\WBth@definition\th@definition
  \let\WBth@remark\th@remark
  \def\th@plain{\WBth@plain\thm@headpunct{。}}%
  \def\th@definition{\WBth@definition\thm@headpunct{。}}%
  \def\th@remark{\WBth@remark\thm@headpunct{。}}%
}
\makeatother

% 无编号小标题（仿 notes3 的 “A Note on ...”）
\newcommand{\notehead}[1]{\par\medskip\noindent{\bfseries #1}\par\smallskip}

% ---------- 表格 ----------
\usepackage{booktabs}
\usepackage{array}
\renewcommand{\arraystretch}{1.25}
\setlength{\tabcolsep}{8pt}

% ---------- 图形 ----------
\usepackage{graphicx}
\usepackage{tikz}
\usepackage{pgfplots}
\pgfplotsset{compat=1.18}
\usetikzlibrary{arrows.meta,positioning,shapes.geometric,calc,patterns,decorations.pathreplacing,decorations.markings}
\usepackage{float}
\usepackage{caption}
\captionsetup{font=small,labelfont=bf,labelsep=quad,skip=4pt}

% 统一图形配色（黑白模式：黑/灰梯度，供所有 TikZ 图复用）
% —— 全局黑白化：原彩色调色板统一改为黑/灰，保证打印与投影均可读 ——
\definecolor{AxisBlue}{RGB}{0,0,0}        % 轴线 / 主曲线：黑
\definecolor{SoftBlue}{RGB}{130,130,130}  % 节点填充：中灰
\definecolor{LightRed}{RGB}{224,224,224}  % 浅填充：浅灰
\definecolor{DeepBlue}{RGB}{0,0,0}        % 描边：黑
\definecolor{MidBlue}{RGB}{70,70,70}      % 次曲线 / 网格：深灰
\tikzset{
  >=Stealth,
  every picture/.append style={font=\small},
  axisline/.style={thick,color=black,-Stealth},
  axislabel/.style={color=black},
  fillsoft/.style={fill=LightRed,fill opacity=0.6,draw=black,thick},
  nodeblue/.style={circle,fill=black!85,draw=black,thick,inner sep=2pt},
  nodewhite/.style={circle,fill=white,draw=black,thick,inner sep=2pt},
  arrblue/.style={->, thick, color=black},
  boxsoft/.style={draw=black, rounded corners=2pt, align=center, inner sep=5pt, fill=black!6},
  % 多曲线 / 多集合区分：以线型 + 标记代替颜色
  curveA/.style={thick,black,solid},
  curveB/.style={thick,black,dashed},
  curveC/.style={thick,black,dotted},
  markA/.style={mark=*,mark size=2.2pt,color=black},
  markB/.style={mark=square*,mark size=2.0pt,color=black},
  markC/.style={mark=triangle*,mark size=2.4pt,color=black},
  % 图内区域 / 点标签（白底避免与图形重叠）
  figlabel/.style={font=\scriptsize,fill=white,inner sep=1pt,outer sep=1pt,text=black},
}
% 图形说明（黑白模式下对图形区域 / 曲线 / 关键点的解释，置于图形下方）
\newcommand{\figexplan}[1]{%
  \par\smallskip\noindent{\small\itshape 图形说明：#1}\par\smallskip}

% ---------- 页眉页脚（书稿格式：左=书名，右=章次与章名）----------
\usepackage{fancyhdr}
% 章标题：居中、大号（书稿章首样式）
\CTEXsetup[format={\centering\LARGE\bfseries}]{chapter}
\CTEXsetup[beforeskip=16pt,afterskip=14pt]{chapter}
\CTEXsetup[format={\large\bfseries}]{section}
\CTEXsetup[format={\normalsize\bfseries}]{subsection}
\pagestyle{fancy}
\fancyhf{}
\providecommand{\booktitle}{微观经济学理论：可行集、偏好与效用}
\providecommand{\lectureno}{第一章}
\providecommand{\lectureshort}{课程讲义}
\fancyhead[L]{\small\itshape \booktitle}
\fancyhead[R]{\small \lectureno\quad \lectureshort}
\fancyfoot[C]{\small \thepage}
\renewcommand{\headrulewidth}{0.4pt}
\fancypagestyle{plain}{\fancyhf{}\fancyfoot[C]{\small \thepage}\renewcommand{\headrulewidth}{0pt}}

% ---------- 正文细节 ----------
\usepackage{setspace}
\setstretch{1.12}
\usepackage{microtype}
\tolerance=900
\emergencystretch=2em
\raggedbottom

% ---------- 超链接与 PDF 书签 ----------
\usepackage{hyperref}
\hypersetup{
  colorlinks=true,
  linkcolor=black,
  urlcolor=black,
  citecolor=black,
  pdfencoding=auto,
  bookmarks=true,
  pdftitle={\booktitle\quad \lectureno：\lectureshort},
  pdfsubject={微观经济学理论书稿}
}
% 章节标题中常见数学符号的安全书签映射（避免 "Improper alphabetic constant"）
\pdfstringdefDisableCommands{%
  \def\leftrightarrow{<->}%
  \def\star{*}%
  \def\S{S}%
  \def\succeq{\text{≽}}%
  \def\succ{\text{≻}}%
  \def\sim{\text{∼}}%
}

% ---------- 常用记号 ----------
\newcommand{\R}{\mathbb{R}}
\newcommand{\N}{\mathbb{N}}
\newcommand{\Z}{\mathbb{Z}}
\newcommand{\Q}{\mathbb{Q}}
\newcommand{\succeqq}{\succeq}
\newcommand{\argmax}{\operatorname*{arg\,max}}
\newcommand{\argmin}{\operatorname*{arg\,min}}
\newcommand{\st}{\text{s.t.}\ }

\setcounter{chapter}{5}
% 缺字回退：主字体缺 ①②③④ 等字形时回退到 Songti SC
\directlua{luaotfload.add_fallback("wbfallback", {"name:Songti SC:mode=node;"})}
\setmainfont{TeX Gyre Termes}[Ligatures=TeX,RawFeature={fallback=wbfallback}]

\begin{document}
\chapter{从需求数据到偏好：可积性、显示偏好与福利评价}
\noindent\textbf{本讲任务。}前五章走的是一条单向的路：从偏好出发，经表示定理得到效用函数，
再经最优化得到需求，最后经对偶与分解得到 Slutsky 矩阵。本章反过来问：\textbf{如果只拿到需求数据，
能退回多远？}这不是数学游戏。任何一项需求估计的实证工作，最终都要回答三个问题：估计出的需求系统
是否符合理论（可检验性）、它是否来自某个偏好（可理性化）、据此算出的福利变化是否有意义（货币度量）。
本章按十个小节展开：
(1) 问题的提法与需求函数的信息含量；
(2) 需求系统的可观测结构与三条加总恒等式；
(3) 对称性在什么维数下才是实质约束；
(4) 路径积分与可积性定理（Hurwicz--Uzawa）；
(5) 三个可算系统的显式恢复（CD、CES、LES）与计量约束检验；
(6) 非参数路径：WARP、SARP、GARP 与 Afriat 定理；
(7) 加总问题与 Gorman 极化形式；
(8) 福利评价：补偿变差、等价变差与消费者剩余；
(9) 一条从数据到福利的实证工作流；
(10) 本章小结（逻辑链、检查清单与常见陷阱）。
全文统一记号：Marshallian 需求 $x(p,w)$、Hicksian 需求 $h(p,\bar u)$、间接效用 $v(p,w)$、
支出函数 $e(p,\bar u)$，Slutsky 矩阵记作 $S(p,w)=D_px(p,w)+\bigl(D_wx(p,w)\bigr)x(p,w)^{\!\top}$。
\tableofcontents
\clearpage
\section{问题的提出：需求数据能告诉我们什么}
% ============================================================
% 第六讲 6.1：反问题的三种提法与需求函数的信息含量
% ============================================================

\subsection{正向链条与它的三个反问题}

前五章建立的是一条单向的推导链。用一张图概括：

\begin{center}
\begin{tikzpicture}[>=Stealth]
  \node[boxsoft] (pref) at (0,0) {偏好 $\succsim$};
  \node[boxsoft] (util) at (4.1,0) {效用 $u(\cdot)$};
  \node[boxsoft] (dem)  at (8.5,0) {需求 $x(p,w)$};
  \node[boxsoft] (dual) at (8.5,-2.7) {支出函数 $e(p,\bar u)$、\\[1pt] Hicksian 需求 $h(p,\bar u)$};
  \draw[arrblue] (pref)--node[figlabel,above,align=center]{表示定理\\[-2pt]{\scriptsize Debreu}} (util);
  \draw[arrblue] (util)--node[figlabel,above,align=center]{最优化\\[-2pt]{\scriptsize Kuhn--Tucker}} (dem);
  \draw[arrblue] (dem)--node[figlabel,right,align=left]{对偶\\[-2pt]{\scriptsize Shephard 引理}} (dual);
  \node[figlabel,below=0.18cm of dual,align=center]
       {Slutsky 矩阵：$S=D_p^2e=\nabla_ph$};
  \draw[thick,dashed,->] (dem.north) to[bend right=26]
       node[figlabel,above,align=center]{反问题：可积性（本章）} (pref.north);
\end{tikzpicture}
\end{center}
\figexplan{实线为前五章建立的正向链条：偏好经表示定理得到效用，经最优化得到需求，再经对偶得到支出函数与 Hicksian 需求（支出函数的 Hessian 即 Slutsky 矩阵，这是全章的枢纽）。上方那条虚线 arc 是本章要回答的反问题：从最右端的 $x(p,w)$ 能否退回最左端的偏好，退回时丢失了什么。注意反问题不必经过效用函数——\S 6.4 的路线是直接把需求积分成支出函数，再由对偶反造偏好。}

反问题可以精确化为三个层次，难度依次上升：

\begin{enumerate}
\item \textbf{函数层面}（可积性）。已知处处光滑的需求函数 $x:\R^n_{++}\times\R_{++}\to\R^n_+$ 的解析式，
      是否存在一个偏好 $\succsim$ 使 $x(p,w)$ 恰为它在预算约束下的需求？
\item \textbf{数据层面}（可理性化）。只观察到有限个价格—数量对 $\{(p^t,x^t)\}_{t=1}^{T}$，
      是否存在一个效用函数与之一致？注意这里连"需求函数"都不知道，只有点。
\item \textbf{福利层面}（货币度量）。若前两层通过，能否用已经恢复出来的 $e(p,\bar u)$
      计算价格变化对消费者的货币影响（补偿变差、等价变差）？
\end{enumerate}

三个层次分别对应本章的 \S 6.2--6.5、\S 6.6、\S 6.8。三个层次的共同前提是一个方法论立场：
\textbf{理论对数据施加了可检验的限制}。如果任何光滑函数都能充当需求函数，需求估计就不存在"设定检验"这一说；
恰恰因为存在限制，计量经济学家才能用数据去拒绝或接受理论。本章的技术内容，本质上就是把这组限制写清楚、
证明它充分、并说明它在什么维数下才有内容。

\subsection{三条基本观察}

在动笔之前必须先弄清楚：反问题\emph{不可能}有唯一的解。原因有三，逐条陈述。

\notehead{观察一：从偏好到需求是多对一的，因此只能恢复到序数水平。}

设 $f:\R\to\R$ 严格递增，则 $f\circ u$ 与 $u$ 刻画同一偏好，从而生成同一个需求函数。
更一般地，任何保持序关系的单调变换都不改变任何可观测的选择。因此反问题的解集若含一个效用函数，
必含它的一整条单调变换轨道。这不是技术缺陷，而是效用理论的既有性质：效用是序数概念，
需求数据也不可能提供比序数更多的信息。\textbf{能期望的最强结论是"恢复到一个单调变换为止"。}

\begin{example}[CD 的两条等价轨道]
取 $u(x)=x_1^{1/2}x_2^{1/2}$ 与 $u'(x)=x_1x_2$。两者生成同一个需求
\[
 x(p,w)=\Bigl(\frac{w}{2p_1},\;\frac{w}{2p_2}\Bigr),
\]
因为 $u'=(u)^2$ 是单调变换。但若以 $u$ 为效用的度量，$e(p,\bar u)=\bar u\,(4p_1p_2)^{1/2}$；
以 $u'$ 为度量，$e'(p,\bar u')=\bar u'\,(p_1p_2)^{1/2}$。两个支出函数不同，然而它们描述的是同一组无差异曲线。
这提示：\textbf{支出函数的值依赖效用标度，但支出函数的两阶导数（即 Slutsky 矩阵）不依赖}——
后者才是需求数据能确定的东西。
\end{example}

\notehead{观察二："序数"不等于"没有内容"，可确定的是曲率而非水平。}

虽然效用的绝对水平不可识别，但无差异曲线的\emph{形状}是唯一的（在同一偏好类内）。
形状在支出函数上的体现是它对价格的二阶导数：$D_p^2e=\nabla_p h$。
这正是可观测的 Slutsky 矩阵。因此本讲的核心结论可以先用一句话预告：
\textbf{需求数据唯一确定偏好的序数结构，其可检验内容集中在 Slutsky 矩阵的对称性与负半定性上。}
而且一旦选定效用标度（例如规定 $e(p,\bar u)=\bar u$ 对某个基准价格向量成立），
支出函数就被完全钉死，福利度量随之确定。

\notehead{观察三：需求函数不能是任意的，否则理论不可检验。}

第三条观察最好用一个具体的失败例子说明。考虑两期观测，两商品，价格向量相同：

\begin{example}[同一价格下选两个不同束]\label{ex:twobundles}
设 $p^1=p^2=(1,1)$，观测到 $x^1=(1,0)$、$x^2=(0,1)$。
若这两期数据来自同一个按预算约束最大化的消费者，则两期的财富必须相同（否则预算集不同，价格相同而"同一价格下两个束都可负担"这句话就不成立）：
$w^1=p^1\cdot x^1=1$，$w^2=p^2\cdot x^2=1$。于是在第 1 期，$x^2$ 也恰好可负担（$p^1\cdot x^2=1\le w^1$），
而消费者选了 $x^1$，说明 $u(x^1)\ge u(x^2)$；对称地，第 2 期说明 $u(x^2)\ge u(x^1)$。
两式相容，并不矛盾——但若把两式都取严格（例如观测到 $x^1$ 时 $x^2$ 严格便宜），矛盾立刻出现。
这就是弱显示偏好公理（WARP）的原型，见 \S 6.6。
\end{example}

这个例子还说明一件事：\textbf{理论限制可以被数据违反}，因此可以被检验。
第 5 章已经证明"任何由偏好生成的需求都满足 Slutsky 对称与负半定"；
本章要做的是反过来证明"满足这两条的需求都能由某个偏好生成"（充分性），
从而把这两条确立为\emph{充分必要}的可理性化条件。定理的两半合起来才构成完整的可检验性陈述：
一半告诉我们理论禁止什么，另一半告诉我们理论允许什么。

\subsection{两条恢复路径与一个警告}

从技术上看，恢复有两条互不相同的路线，本章两条都走：

\begin{center}
\footnotesize
\begin{tabular}{l>{\raggedright\arraybackslash}p{3.7cm}>{\raggedright\arraybackslash}p{5.4cm}l}
\toprule
路径 & 所需信息 & 核心工具 & 本章位置\\
\midrule
微分路径（参数方法） & 需求函数解析式、可微 & 路径积分、Frobenius 条件 & \S 6.2--6.5\\
非参数路径 & 有限观测点 $\{(p^t,x^t)\}$ & 显示偏好、Afriat 定理、线性规划 & \S 6.6\\
\bottomrule
\end{tabular}
\end{center}
\figexplan{两条路径的分工：微分路径假定知道整个函数，可以得到闭式的 $e$ 与 $h$，代价是需要可微性与具体的函数形式；非参数路径只用有限数据点，不需要任何函数形式，但结论是"存在一个效用函数"而非"这个效用函数"，且不能直接给出 $e$ 的解析式。两条路径在理论上相容：若数据来自一个可积的需求系统，则它必满足 GARP。}

\begin{remark}[一个必须先说清的警告：局部与全局]
可积性定理是\emph{局部}陈述：在价格空间的某个凸（或星形）邻域上，一组满足条件的偏导数可以积出一个支出函数。
把它推广到整个 $\R^n_{++}$ 需要额外的技术条件（如解曲线不逃逸、凸性在整体上成立）。
本讲按教科书传统在"相关区域"内陈述定理，并把技术边界显式标出，避免出现"定理说了全局而证明只做了局部"这类常见疏漏。
第 5 章对"相关区间"的处理（$\bar u\in(u(0),\sup u)$）出于完全相同的谨慎。
\end{remark}

\begin{remark}[术语小注]
文献中 \emph{integrability}（可积性）指"需求函数可积出支出函数"；\emph{rationalizability}（可理性化）指
"数据可由某个偏好生成"；\emph{integrability conditions} 与 \emph{Slutsky conditions} 常混用，
但严格说来前者包含后者加上可微性与边界条件。McKenzie（1957）曾指出一个方法论要点：
Slutsky 条件可以不借助任何效用概念、直接在需求函数上表述并赋予行为解释，
这正是后来显示偏好纲领的技术先声。本讲采纳这一立场——先写条件，再谈偏好。
\end{remark}

\section{需求系统的可观测结构：三个恒等式与 Slutsky 矩阵}
% ============================================================
% 第六讲 6.2：正则性假设、Engel/Cournot/齐次恒等式、S 的弹性形式
% ============================================================

\subsection{正则性假设}

反问题要能提得清楚，先要给"需求函数"划一个可操作的定义域与光滑性台阶。
以下四条贯穿全章，除特别声明外不再重复。

\begin{condition}[正则性]\label{cond:regular}
需求函数 $x:\R^n_{++}\times\R_{++}\to\R^n_+$ 满足：
\begin{enumerate}
\item[\textbf{(A1)}] \textbf{光滑性}：$x$ 连续可微，且在各分量上足够多次可微（用到几阶就假设几阶）；
\item[\textbf{(A2)}] \textbf{零次齐次}：$x(tp,tw)=x(p,w)$，$\forall t>0$；
\item[\textbf{(A3)}] \textbf{Walras 律}：$p\cdot x(p,w)=w$，$\forall(p,w)$；
\item[\textbf{(A4)}] \textbf{内点性}：$x(p,w)\gg0$，从而不涉及角点与库恩—塔克乘子的符号讨论。
\end{enumerate}
\end{condition}

\begin{remark}[每条假设在做什么]
(A1) 是为了能写导数、能积分，是纯技术条件。(A2) 是"没有货币幻觉"：价格与财富同比例变动，
可行集不变，选择也不变。(A3) 说预算被花光，其理论依据是局部非饱和（第 5 章 \S 5.1）；
注意 (A3) 是\emph{等式}，它把 $n$ 维需求压缩到 $n-1$ 维，这一点将在下文反复用到。
(A4) 排除了角点，代价是丢掉"某些商品不消费"这一大类现实情形；\S 6.6 的非参数方法不需要 (A4)，
这也是它在实证中更受欢迎的原因之一。
\end{remark}

\subsection{三条加总恒等式}

把 (A2)、(A3) 对参数求导，立刻得到三组恒等式。它们在实证中不是"可选的美化"，而是任何需求系统
都必须逐点满足的会计约束；估计时若不施加，参数估计就与预算约束自相矛盾。

\begin{proposition}[Engel 加总与 Cournot 加总]\label{prop:aggr-identities}
在 (A1)--(A3) 下，对一切 $(p,w)$：
\begin{align}
\text{（Engel 加总）}\qquad & \sum_{i=1}^{n}p_i\,\frac{\partial x_i(p,w)}{\partial w}=1,
\label{eq:engel}\\[2pt]
\text{（Cournot 加总）}\qquad & x_j(p,w)+\sum_{i=1}^{n}p_i\,\frac{\partial x_i(p,w)}{\partial p_j}=0,
\qquad j=1,\ldots,n,
\label{eq:cournot}\\[2pt]
\text{（齐次性）}\qquad & \sum_{j=1}^{n}p_j\,\frac{\partial x_i(p,w)}{\partial p_j}
 +w\,\frac{\partial x_i(p,w)}{\partial w}=0,\qquad i=1,\ldots,n.
\label{eq:homog}
\end{align}
\end{proposition}

\begin{proof}
\eqref{eq:engel}：把 $p\cdot x(p,w)=w$ 对 $w$ 求导，$p$ 视为常数，得 $\sum_i p_i\partial x_i/\partial w=1$。
\eqref{eq:cournot}：把同一恒等式对 $p_j$ 求导，注意左端的 $p$ 也依赖 $p_j$：
$\partial\bigl(\sum_i p_ix_i\bigr)/\partial p_j=x_j+\sum_i p_i\,\partial x_i/\partial p_j=0$，右端 $w$ 与 $p_j$ 无关。
\eqref{eq:homog}：把 $x_i(tp,tw)=x_i(p,w)$ 对 $t$ 求导后在 $t=1$ 处取值，即得 Euler 恒等式。
\end{proof}

\begin{remark}[三条恒等式的经济读法]
Engel 加总说的是"新增一元财富必须全部花掉"：各商品的财富导数按价格加权后恰为 1，
因此不可能所有商品都是劣等品（否则加权和 $\le0$）。
Cournot 加总说的是"第 $j$ 种商品涨价，若要保持原来的消费束，支出恰好增加 $x_j$"——
价格上升的全部直接效应被数量完全吸收。
齐次性说的是"价格与财富同比例上升时，各商品需求的净变化为零"：价格效应与财富效应精确抵消。
三条恒等式合起来把 $n\times n$ 的价格导数矩阵与 $n$ 维财富导数绑在一起，
自由参数的个数因此从 $n^2+n$ 降到 $n^2-n$（这一点在 \S 6.5 的计量约束计数中会再出现）。
\end{remark}

\subsection{Slutsky 矩阵：定义与自动成立的两条性质}

\begin{definition}[Slutsky 矩阵]\label{def:slutsky}
在 (A1)--(A3) 下定义
\[
 \boxed{\;S(p,w)\;=\;D_px(p,w)+\bigl(D_wx(p,w)\bigr)\,x(p,w)^{\!\top}\;}
\]
其第 $(i,j)$ 元为
\[
 S_{ij}(p,w)=\frac{\partial x_i(p,w)}{\partial p_j}
 +x_j(p,w)\,\frac{\partial x_i(p,w)}{\partial w}.
\]
\end{definition}

注意这个定义的\textbf{全部右端项都是可观测的}：价格导数、财富导数与需求量本身。
这正是 \S 6.3 之后所有检验得以进行的前提——Slutsky 矩阵不是从偏好算出来的理论对象，
而是可以从数据估出来的统计对象。

\begin{proposition}[两条自动成立的性质]\label{prop:S-auto}
在 (A1)--(A3) 下，对一切 $(p,w)$：
\begin{enumerate}
\item $S(p,w)\,p=0$（行和为零，按价格加权）；
\item $p^{\top}S(p,w)=0$（列和为零，按价格加权），从而 $\operatorname{rank}S\le n-1$。
\end{enumerate}
\end{proposition}

\begin{proof}
(1) 用 \eqref{eq:homog} 的矩阵形式 $D_px\cdot p+w\,D_wx=0$：
\[
 S\,p=D_px\cdot p+\bigl(D_wx\bigr)\bigl(x^{\top}p\bigr)
 =D_px\cdot p+w\,D_wx=0,
\]
其中第二个等号用了 Walras 律 $x^{\top}p=w$。

(2) 用 \eqref{eq:cournot} 的矩阵形式 $p^{\top}D_px=-x^{\top}$，以及 \eqref{eq:engel} 的 $p^{\top}D_wx=1$：
\[
 p^{\top}S=p^{\top}D_px+\bigl(p^{\top}D_wx\bigr)x^{\top}
 =-x^{\top}+1\cdot x^{\top}=0.
\]
由于 $p\gg0$ 非零，$S$ 的零空间至少含 $\operatorname{span}\{p\}$，故秩不超过 $n-1$。
\end{proof}

\begin{remark}[不要把自动性质当成检验]
命题~\ref{prop:S-auto} 常被误引为"Slutsky 理论的经验验证"。这是误解：
行和为零、列和为零是齐次性与 Walras 律的直接推论，而这两条在估计时通常被\emph{直接施加}（作为参数约束）。
真正有内容、可以拒绝理论的是\textbf{对称性}与\textbf{负半定性}——它们不是会计恒等式，
而是偏好最大化的实质后果。\S 6.3 与 \S 6.5 分别从理论与数值两侧说明这一点。
\end{remark}

\subsection{弹性形式：实证中的标准写法}

计量工作几乎从不直接报告 $S$ 的原始元，而是报告弹性。把 $S$ 翻译成弹性形式，
既便于跨商品、跨样本比较，也让约束的形式变得一目了然。

记支出份额 $s_i=p_ix_i/w$，价格弹性 $\varepsilon_{ij}=(\partial x_i/\partial p_j)(p_j/x_i)$，
财富弹性 $\eta_i=(\partial x_i/\partial w)(w/x_i)$。由定义~\ref{def:slutsky} 直接展开：

\begin{proposition}[Slutsky 矩阵的弹性形式]\label{prop:S-elasticity}
在 (A1)--(A3) 且 $x_i>0$ 下，定义 Slutsky 弹性 $\sigma_{ij}=\varepsilon_{ij}+s_j\eta_i$，则
\[
 S_{ij}=\frac{x_i}{p_j}\,\sigma_{ij},
 \qquad\text{等价地}\qquad
 \Sigma_{ij}:=\frac{p_ip_j}{w}S_{ij}=s_i\,\sigma_{ij}.
\]
并且：
\begin{enumerate}
\item \textbf{对称性} $\iff$ $s_i\sigma_{ij}=s_j\sigma_{ji}$，$\forall i,j$；
\item \textbf{负半定性} $\iff$ 矩阵 $\Sigma=(\Sigma_{ij})$ 负半定；
\item \textbf{行和为零} $\iff$ $\sum_j s_j\sigma_{ij}=0$，$\forall i$。
\end{enumerate}
\end{proposition}

\begin{proof}
由 $S_{ij}=\partial x_i/\partial p_j+x_j\,\partial x_i/\partial w$，在第一项提出 $x_i/p_j$、
第二项提出 $x_i/w$：
\[
 S_{ij}=\frac{x_i}{p_j}\cdot\frac{p_j}{x_i}\frac{\partial x_i}{\partial p_j}
 +x_j\cdot\frac{x_i}{w}\cdot\frac{w}{x_i}\frac{\partial x_i}{\partial w}
 =\frac{x_i}{p_j}\,\varepsilon_{ij}+\frac{x_ix_j}{w}\,\eta_i
 =\frac{x_i}{p_j}\bigl(\varepsilon_{ij}+s_j\eta_i\bigr),
\]
末一步用了 $x_j/w=s_j/p_j$。两端乘 $p_ip_j/w$ 得 $\Sigma_{ij}=s_i\sigma_{ij}$。
(1)(2)(3) 分别是"对称""二次型符号""行和"在变量替换 $S\mapsto\Sigma=\operatorname{diag}(p)S\operatorname{diag}(p)/w$ 下的不变性：
该变换是同余变换（合同变换），保持对称性与惯性，故负半定性不变。
\end{proof}

\begin{remark}[为什么用 $\Sigma$ 而不是 $S$ 做检验]
$\Sigma$ 的元素是无量纲的（份额乘弹性），可以直接跨商品比较，也可以在同一个矩阵里同时读替代与互补。
更重要的是：$\Sigma$ 的对称性条件 $s_i\sigma_{ij}=s_j\sigma_{ji}$ 让"谁替代谁"的非对称性变得可见——
原始形式 $S_{ij}=S_{ji}$ 掩盖了份额差异，而份额差异恰恰是实证中拒绝对称性的主要来源之一。
\end{remark}

把本节的结果整理成一张清单，作为后文的索引：

\begin{center}
\begin{tabular}{llll}
\toprule
条件 & 数学表述 & 来源 & 是否实质约束\\
\midrule
齐次性 & $D_px\cdot p+wD_wx=0$ & 无货币幻觉 & 否（可施加）\\
Walras 律 & $p\cdot x=w$ & 局部非饱和 & 否（可施加）\\
行和为零 & $Sp=0$ & 上两条的推论 & 否\\
列和为零 & $p^{\top}S=0$ & 上两条的推论 & 否\\
\textbf{对称性} & $S=S^{\top}$ & 支出函数二阶导对称 & \textbf{是}（$n\ge3$）\\
\textbf{负半定性} & $v^{\top}Sv\le0$ & $e$ 对 $p$ 凹 & \textbf{是}（一切 $n\ge2$）\\
\bottomrule
\end{tabular}
\end{center}
\figexplan{可积性条件的清单。前四行由正则性假设直接给出，估计时作为参数约束施加即可；后两行才是理论对数据的真正限制，也是本章的检验对象。注意对称性在 $n=2$ 时自动成立（\S 6.3 命题~\ref{prop:n2-auto}），因此两商品情形下唯一可检验的内容是曲率。}

\section{对称性：什么时候是实质约束}
% ============================================================
% 第六讲 6.3：n=2 自动对称、n≥3 实质约束、曲率成为唯一内容
% ============================================================

\subsection{必要性：从偏好到对称}

先把第 5 章已证的一半结论固定下来，作为本章的出发点。

\begin{theorem}[Slutsky 条件的必要性]\label{thm:slutsky-necessity}
设偏好 $\succsim$ 连续、严格凸、局部非饱和，$u$ 为其连续表示，需求 $x(p,w)$ 在 (A1) 下可微。
则对一切 $(p,w)$，$S(p,w)$ 对称且负半定。
\end{theorem}

\begin{proof}
由第 5 章 \S 5.7 的对偶恒等式，$x(p,w)=h\bigl(p,v(p,w)\bigr)$。在 Hicksian 需求上取导数：
\[
 \frac{\partial x_i}{\partial p_j}
 =\frac{\partial h_i}{\partial p_j}
 +\frac{\partial h_i}{\partial \bar u}\,\frac{\partial v}{\partial p_j}
 =\frac{\partial h_i}{\partial p_j}
 -\frac{\partial h_i}{\partial \bar u}\,\lambda\,x_j,
\]
末一步用了 Roy identity $\partial v/\partial p_j=-\lambda x_j$（$\lambda=\partial v/\partial w$）。
再由 EMP 的包络结果 $\partial e/\partial\bar u=\mu$ 与 Hicksian 需求对效用的导数关系
$\partial h_i/\partial\bar u=\partial^2e/(\partial p_i\partial\bar u)$，有
$\partial h_i/\partial\bar u=\mu\,\partial x_i/\partial w$（在 $\bar u=v(p,w)$ 处取值；这可由对偶恒等式
$h(p,v(p,w))=x(p,w)$ 对 $w$ 求导、并用 $\partial v/\partial w=\lambda$ 得到）。
代入上式并用 $\lambda\mu=1$（第 5 章 \S 5.7 的乘子互逆关系）：
\[
 \frac{\partial x_i}{\partial p_j}=\frac{\partial h_i}{\partial p_j}-x_j\frac{\partial x_i}{\partial w},
 \qquad\text{即}\qquad
 S_{ij}=\frac{\partial h_i}{\partial p_j}.
\]
于是 $S=D_ph=\nabla^2_{pp}e$。由 Young 定理，Hessian 对称；由 $e$ 对 $p$ 凹，Hessian 负半定。
\end{proof}

\begin{remark}[这一步的意义]
定理~\ref{thm:slutsky-necessity} 把"可观测的 $S$"与"理论上的 $\nabla^2_{pp}e$"直接等同。
这是全章的技术枢纽：对称性与负半定性之所以能作为检验，正因为它们恰好是支出函数 Hessian 的两条性质。
换言之，$S$ 就是 $e$ 的曲率在需求一侧的影子。
\end{remark}

\subsection{两商品情形：对称性自动成立}

初学者常把"$S$ 对称"当作一条需要单独验证的独立条件。在两商品世界里，这条条件\emph{没有任何内容}——
它被齐次性与 Walras 律自动蕴含。

\begin{proposition}[$n=2$ 时对称性自动成立]\label{prop:n2-auto}
设 $n=2$，$x=(x_1,x_2)$ 满足 (A1)--(A3)，$p\gg0$，$x_1>0$。则 $S_{12}=S_{21}$，即 $S$ 必对称。
\end{proposition}

\begin{proof}
定义\emph{预算份额函数}
\[
 g(r):=\frac{p_1x_1(p,w)}{w},\qquad r:=\frac{p_2}{p_1}.
\]
由零次齐次性 (A2)，$x_1(tp,tw)=x_1(p,w)$，故 $p_1x_1/w$ 只依赖比值 $r$：$g$ 确实是单变量函数。
由此 $x_1=wg(r)/p_1$，再由 Walras 律 (A3) 解出
\[
 x_2=\frac{w-p_1x_1}{p_2}=\frac{w\bigl(1-g(r)\bigr)}{p_2}.
\]
于是整个需求系统被\emph{一个}单变量函数 $g$ 完全刻画——这正是"$n=2$ 时自由度只剩一个"的含义。
逐项求导（$r=p_2/p_1$，故 $\partial r/\partial p_1=-p_2/p_1^2$，$\partial r/\partial p_2=1/p_1$）：
\begin{align*}
\frac{\partial x_1}{\partial p_1}&=-\frac{wg}{p_1^2}-\frac{wg'p_2}{p_1^3},&
\frac{\partial x_1}{\partial p_2}&=\frac{wg'}{p_1^2},&
\frac{\partial x_1}{\partial w}&=\frac{g}{p_1},\\
\frac{\partial x_2}{\partial p_1}&=\frac{wg'}{p_1^2},&
\frac{\partial x_2}{\partial p_2}&=-\frac{wg'}{p_1p_2}-\frac{w(1-g)}{p_2^2},&
\frac{\partial x_2}{\partial w}&=\frac{1-g}{p_2}.
\end{align*}
代入定义~\ref{def:slutsky}：
\begin{align*}
S_{12}&=\frac{\partial x_1}{\partial p_2}+\frac{\partial x_1}{\partial w}x_2
 =\frac{wg'}{p_1^2}+\frac{g\,(1-g)w}{p_1p_2},\\
S_{21}&=\frac{\partial x_2}{\partial p_1}+\frac{\partial x_2}{\partial w}x_1
 =\frac{wg'}{p_1^2}+\frac{(1-g)\,g\,w}{p_1p_2}.
\end{align*}
两式右端逐项相同，故 $S_{12}=S_{21}$。
\end{proof}

\begin{remark}[$n=2$ 时的显式结构]
在命题~\ref{prop:n2-auto} 的记号下还可算出 $S_{11}$，于是整个矩阵可写成秩一形式
\[
 S=\alpha\begin{pmatrix}1\\[1pt] -p_1/p_2\end{pmatrix}
   \begin{pmatrix}1 & -p_1/p_2\end{pmatrix},
 \qquad
 \alpha=S_{11}=-\frac{wg}{p_1^2}-\frac{wg'p_2}{p_1^3}+\frac{wg^2}{p_1^2}.
\]
$\operatorname{rank}S=1=n-1$，与命题~\ref{prop:S-auto} 的上界吻合。更重要的是：
\textbf{负半定性等价于单个标量条件 $\alpha\le0$}，而这个 $\alpha$ 恰是补偿需求曲线的斜率
$\partial h_1/\partial p_1$。所以在两商品世界里，可积性的全部实质内容就是"补偿需求曲线向下倾斜"。
\end{remark}

\subsection{曲率并不自动：一个两商品反例}

既然对称性在两商品时是免费的，是否意味着两商品的任何"看起来合理"的需求系统都能由偏好生成？
不是。命题~\ref{prop:n2-auto} 只给出对称，负半定性仍可失败。

\begin{example}[对称但不可积的两商品系统]\label{ex:n2-curvature}
取
\[
 x_1(p,w)=\frac{w}{p_1}\cdot\frac{1}{1+(p_2/p_1)^2},\qquad
 x_2(p,w)=\frac{w-p_1x_1(p,w)}{p_2}.
\]
（即命题~\ref{prop:n2-auto} 中取 $g(r)=1/(1+r^2)$。）
在 $p=(2,3)$、$w=50$ 处：$g(1.5)=1/3.25\approx0.3077$，
\[
 x=(7.6923,\;11.5385),\qquad p\cdot x=50,
\]
Slutsky 矩阵为
\[
 S=\begin{pmatrix}
 \phantom{-}2.6627 & -1.7751\\
 -1.7751 & \phantom{-}1.1834
 \end{pmatrix}.
\]
它对称（命题~\ref{prop:n2-auto} 已保证），行和为零（$2.6627\times2-1.7751\times3\approx0$），
但 $S_{11}=2.6627>0$。
\end{example}

\begin{remark}[这个反例说明了什么]
$S_{11}>0$ 意味着：\emph{补偿}（即保持效用不变）地提高商品 1 的价格，商品 1 的需求反而上升。
这与"替代效应必为负"直接冲突，因而违反补偿需求定律，不可能来自任何凸偏好。
值得注意的是这个系统\textbf{处处满足}齐次性、Walras 律乃至对称性——它在会计上完美无缺，
在经济学上却不可能。结论：\textbf{在两商品情形下，可积性的实质内容是曲率，而不是对称}。
\end{remark}

\subsection{三商品及以上：对称性成为真正的限制}

\begin{proposition}[$n\ge3$ 时对称性是实质约束]\label{prop:n3-substantive}
存在满足 (A1)--(A3)、且对称部分负半定的需求系统，其 $S$ 不对称，因而不存在生成它的偏好。
\end{proposition}

\begin{proof}[构造性证明]
取 $n=3$、$\varepsilon>0$，定义
\[
 x_1(p,w)=\frac{w}{3p_1}\Bigl(1+\varepsilon\frac{p_2}{p_1}\Bigr),\qquad
 x_2(p,w)=\frac{w}{3p_2},\qquad
 x_3(p,w)=\frac{w-p_1x_1-p_2x_2}{p_3}.
\]
(A2) 显然（每项对 $(p,w)$ 零次齐次）；(A3) 由 $x_3$ 的定义直接成立；(A1) 成立。
取 $\varepsilon=0.35$、$p=(2,3,5)$、$w=60$，得
\[
 x=(15.25,\;6.6667,\;1.9),\qquad p\cdot x=60,
\]
\[
 S=\begin{pmatrix}
 -6.3740 & \phantom{-}3.4444 & \phantom{-}0.4829\\
 \phantom{-}1.6944 & -1.4815 & \phantom{-}0.2111\\
 \phantom{-}1.5329 & -0.4889 & -0.3198
 \end{pmatrix},
\qquad
 S_{12}=3.4444\ne1.6944=S_{21}.
\]
故 $S$ 不对称，由定理~\ref{thm:slutsky-necessity} 的必要性，不存在生成 $x$ 的偏好。
\end{proof}

\begin{remark}[两个条件互相独立]
上例还揭示了一个容易被忽略的事实：该 $S$ 的\emph{对称部分} $(S+S^{\top})/2$ 的特征值为
$(-7.6077,\,-0.5676,\,0)$，全部非正，因此二次型 $v^{\top}Sv\le0$ 成立。
也就是说，这个系统\emph{满足}补偿需求定律，却\emph{违反}对称性。
可见"负半定"与"对称"是两条逻辑上独立的条件：
前者约束二次型的符号（可观测为补偿需求的自价格效应），后者约束矩阵的型（可观测为交叉价格效应的可积性）。
一个需求系统可以同时满足会计恒等式、满足曲率，却仍然不可积——这正是对称性检验不可替代的原因。
\end{remark}

\begin{remark}[维数与约束数]
把三节的结论并排放，可以看到一个颇有意味的模式：
\begin{center}
\begin{tabular}{llll}
\toprule
维数 & $S$ 的自由元 & 对称性约束数 & 实质约束的载体\\
\midrule
$n=2$ & 1（行和为零后） & 0 & 仅曲率 $\alpha\le0$\\
$n=3$ & 3 & 1 & 对称（1 条）$+$ 曲率\\
$n$ & $n(n-1)/2$ & $n(n-1)/2-(n-1)$ & 对称 $+$ 曲率\\
\bottomrule
\end{tabular}
\end{center}
\figexplan{$S$ 满足 $Sp=0$ 后自由元为 $(n-1)^2$，其中对称性留下 $n(n-1)/2$ 个独立元、反对称部分有 $n(n-1)/2-(n-1)$ 个独立元（后者必须为零才是可积的）。$n=2$ 时反对称部分的维数为 $1-1=0$，故对称性自动成立；$n\ge3$ 时反对称部分非空，对称性才成为可拒绝的约束。这张表也是实证中选择商品数目的一个隐性约束：商品划得越细，理论的可检验内容越多，但估计所需的样本与价格变异也越大。}
\end{remark}

\section{路径积分与可积性定理}
% ============================================================
% 第六讲 6.4：Pfaff 系统、Frobenius 条件、Hurwicz–Uzawa 定理
% ============================================================

\subsection{恢复的对象应是支出函数，而不是效用函数}

先解决一个方向性问题。若试图直接恢复效用函数 $u$，会遇到两个不可逾越的障碍：
其一，$u$ 的取值本身不可识别（观察一的单调变换）；其二，需求只提供\emph{边际}信息，
而 $u$ 的恢复需要把边际信息积分成水平——积分的常数无法由数据确定。

正确的恢复对象是支出函数，理由有三：
\begin{enumerate}
\item 需求与支出函数之间是\emph{梯度关系} $\nabla_pe=h$（Shephard 引理），这是一个可以直接积分的方程；
\item 需求与效用之间没有这样的梯度关系：$x$ 只通过一阶条件与 $\nabla u$ 的\emph{比值}相连，
      比例因子 $\lambda$ 随 $(p,w)$ 变化，无法直接积分；
\item 一旦得到 $e$，效用函数可由对偶重新构造（\S 6.4.5），且这一构造只到单调变换为止，
      正好匹配数据的信息含量。
\end{enumerate}

\subsection{把问题写成偏微分方程组}

设待求的支出函数为 $e(p,\bar u)$，其中 $\bar u$ 作为参数固定。Shephard 引理要求
\[
 \frac{\partial e(p,\bar u)}{\partial p_i}=h_i(p,\bar u),\qquad i=1,\ldots,n.
\]
右端的 Hicksian 需求不可观测，但由对偶恒等式 $h\bigl(p,\bar u\bigr)=x\bigl(p,e(p,\bar u)\bigr)$，
可以把它换成可观测的 Marshallian 需求在\emph{未知支出水平}处的取值：
\begin{equation}\label{eq:pfaff}
 \boxed{\;\frac{\partial e(p,\bar u)}{\partial p_i}
 =x_i\bigl(p,\,e(p,\bar u)\bigr),\qquad i=1,\ldots,n.\;}
\end{equation}
这是一个关于 $e$ 的一阶非线性偏微分方程组（Pfaff 系统）；$\bar u$ 只以参数身份出现，
因此每固定一个 $\bar u$ 就得到一组方程。初始条件取在某个基准价格 $p^0$ 上：
$e(p^0,\bar u)=e^0$，其中 $e^0$ 的选取等价于选定效用的计量单位（见 \S 6.4.5）。

\subsection{可积条件就是对称性}

方程组 \eqref{eq:pfaff} 有解的必要条件是混合偏导可交换。把右端记为 $f_i(p,e)=x_i(p,e)$，
则 $\partial e/\partial p_i=f_i$ 要求
\[
 \frac{\partial f_i}{\partial p_j}+\frac{\partial f_i}{\partial e}\,f_j
 =\frac{\partial f_j}{\partial p_i}+\frac{\partial f_j}{\partial e}\,f_i
 \qquad\text{（Frobenius 条件）}.
\]
逐项计算左端：
\[
 \frac{\partial f_i}{\partial p_j}+\frac{\partial f_i}{\partial e}f_j
 =\frac{\partial x_i}{\partial p_j}+\frac{\partial x_i}{\partial w}\,x_j
 =S_{ij}(p,e).
\]
右端同理为 $S_{ji}(p,e)$。于是得到本节的核心结论：

\begin{proposition}[可积条件 $=$ 对称性]\label{prop:frobenius}
Pfaff 系统 \eqref{eq:pfaff} 在 $(p,\bar u)$ 的邻域内可解（存在 $C^2$ 解 $e$）的必要条件是
$S(p,e(p,\bar u))$ 对称；在 $S$ 连续可微且定义域单连通时，该条件也是充分的。
\end{proposition}

\begin{proof}[证明要点]
必要性即上面的混合偏导计算。充分性用 Frobenius 定理：若系数矩阵对称且光滑，
则方程组的解族存在且依赖一个常数（初值）。单连通性保证任意两条连接 $p^0$ 与 $p$ 的路径给出同一解，
即积分与路径无关；这正是 Poincaré 引理在"系数含未知函数"情形下的推广。
\end{proof}

\begin{remark}[几何读法：路径无关]
图~\ref{fig:path} 给出二维价格空间中的图像。从 $p^0$ 到 $p^1$ 有无数条路径，
沿每条路径把"需求 $\times$ 价格增量"累加，得到若干个数值。若 $S$ 对称，这些数值全部相等，
线积分因此定义了一个函数 $e$；若 $S$ 不对称，沿不同路径累加会得到不同结果，
"支出函数"就不存在——同一组价格与效用会对应多个"最小支出"，这在逻辑上不可能。
所以对称性不是一条外加的公理，而是"支出函数存在"这一句话的微分形式。
\end{remark}

\begin{figure}[H]
\centering
\begin{tikzpicture}[scale=1.0]
  \draw[axisline] (0,0)--(6.4,0) node[below]{$p_1$};
  \draw[axisline] (0,0)--(0,4.9) node[left]{$p_2$};
  \coordinate (P0) at (1.1,1.0);
  \coordinate (P1) at (5.2,3.9);
  \draw[fillsoft] (1.1,1.0)--(5.2,1.0)--(5.2,3.9)--(1.1,3.9)--cycle;
  % 路径 A：先横后竖
  \draw[curveA] (P0)--(5.2,1.0)--(P1);
  % 路径 B：先竖后横
  \draw[curveB] (P0)--(1.1,3.9)--(P1);
  % 路径 C：直线
  \draw[curveC] (P0)--(P1);
  \node[nodewhite] at (P0) {};
  \node[figlabel,below left] at (P0) {$p^0$};
  \node[nodewhite] at (P1) {};
  \node[figlabel,above right] at (P1) {$p^1$};
  \node[figlabel] at (3.35,0.70) {路径 A};
  \node[figlabel,rotate=90] at (0.72,2.55) {路径 B};
  \node[figlabel] at (3.55,2.02) {路径 C};
\end{tikzpicture}
\caption{线积分的路径无关性}
\figexplan{价格空间中的三条路径：A（实线，先调 $p_1$ 再调 $p_2$）、B（虚线，顺序相反）、C（点线，沿直线同时调整）。支出函数的变化量定义为沿路径对 $\sum_i h_i\,\mathrm dp_i$ 的线积分。对称性 $S_{ij}=S_{ji}$ 保证 A、B、C 三条路径给出同一数值，等价于矩形回路上的环路积分为零（图中阴影区域即该回路）；若 $S$ 不对称，绕回路一周的积分不为零，支出函数无法定义。}
\label{fig:path}
\end{figure}

\subsection{Hurwicz--Uzawa 可积性定理}

\begin{theorem}[可积性定理]\label{thm:integrability}
设 $x$ 满足 (A1)--(A3)，且在价格空间的单连通开集 $\mathcal P$ 上
\begin{enumerate}
\item[(i)] $S(p,w)$ 对称；
\item[(ii)] $S(p,w)$ 负半定。
\end{enumerate}
则存在函数 $e:\mathcal P\times(u(0),\infty)\to\R_{++}$ 与一个连续、严格凸、局部非饱和的偏好 $\succsim$，使得：
\begin{enumerate}
\item $e(\cdot,\bar u)$ 对 $p$ 一次齐次、凹、非减，且 $\nabla_pe=h$；
\item $x(p,w)$ 恰为 $\succsim$ 在预算集 $\{x:p\cdot x\le w\}$ 上的需求；
\item $e$ 由 $(x,\bar u)$ 唯一确定（在选定效用标度之后）。
\end{enumerate}
\end{theorem}

\begin{proof}[证明纲要]
分三步。技术细节（可微性阶数、边界行为、解曲线的整体存在性）从略，只保留构造的逻辑骨架。

\textbf{第一步：构造 $e$。}固定 $\bar u$，在 $\mathcal P$ 上取定基准点 $p^0$ 与初值 $e^0>0$，
沿任一条从 $p^0$ 到 $p$ 的路径 $\Gamma$ 求解 \eqref{eq:pfaff}（等价于沿 $\Gamma$ 积分）。
由命题~\ref{prop:frobenius}，条件 (i) 保证解与路径无关，故得到良定义的 $C^2$ 函数 $e(p,\bar u)$。

\textbf{第二步：验证 $e$ 是支出函数。}
\emph{一次齐次}：$x$ 零次齐次意味着 \eqref{eq:pfaff} 在 $p\mapsto tp$ 下保持形式，
解的唯一性给出 $e(tp,\bar u)=t\,e(p,\bar u)$。
\emph{对 $p$ 非减}：$\partial e/\partial p_i=x_i\ge0$。
\emph{凹性}：对方程 \eqref{eq:pfaff} 再求一次偏导，
\[
 \frac{\partial^2e}{\partial p_i\partial p_j}
 =\frac{\partial x_i}{\partial p_j}+\frac{\partial x_i}{\partial w}\frac{\partial e}{\partial p_j}
 =S_{ij},
\]
即 $D_p^2e=S$。条件 (ii) 因此恰好等价于 $e$ 对 $p$ 凹。
\emph{对 $\bar u$ 单调}：由构造中 $\bar u$ 的排序保证（技术细节见 Hurwicz--Uzawa 原文）。

\textbf{第三步：由 $e$ 反造偏好。}定义间接效用为 $e$ 对 $\bar u$ 的广义逆：
\[
 V(p,w):=\max\{\bar u:\;e(p,\bar u)\le w\}.
\]
由 $e$ 对 $\bar u$ 严格递增且连续，$V$ 良定义，且对 $(p,w)$ 零次齐次、拟凸。
再由标准的对偶构造定义直接效用
\[
 u(x):=\min_{p\gg0}\bigl\{V(p,\,p\cdot x)\bigr\},
\]
即在所有可能价格下"最不利"的间接效用。可以验证 $u$ 连续、拟凹，且其生成的 Marshallian 需求恰为 $x$：
由 $V$ 的定义，$x(p,w)$ 满足 $e\bigl(p,V(p,w)\bigr)=w$，再由 $\nabla_pe=x(p,e)$ 与 Shephard 引理，
$x(p,w)=h\bigl(p,V(p,w)\bigr)$；即 $x$ 是支出最小化在效用水平 $V(p,w)$ 处的解，
由对偶恒等式它同时是效用最大化的解。
\end{proof}

\begin{remark}[定理的地位]
定理~\ref{thm:integrability} 与定理~\ref{thm:slutsky-necessity} 合起来构成一个\emph{充分必要}刻画：
\[
 x\ \text{可由偏好生成}
 \iff
 \begin{cases}
 \text{(A1)--(A3)},\\
 S\ \text{对称},\\
 S\ \text{负半定}.
 \end{cases}
\]
这是消费者理论中最重要的"元定理"：它把"理性"这一不可直接观测的假设，
转换成一组可以施加在估计结果上的、有限维的、可拒绝的约束。
没有它，需求估计就只是曲线拟合；有了它，实证才有检验理论的资格。
\end{remark}

\begin{remark}[局部与全局的分界]
定理的陈述限定在单连通开集 $\mathcal P$ 上，这是 Frobenius 定理的自然适用范围。
若要 $e$ 定义在整个 $\R^n_{++}$ 上，需附加条件：解曲线不逃出定义域、$S$ 在整个域上负半定、
以及边界行为（$p_i\to0$ 时 $x_i$ 的极限）不破坏连续性。
实证中这些条件极少被检验，原因很实际：样本只覆盖价格空间的一小块，全局可积性无法从局部数据识别。
因此\textbf{"可积"在实证中始终是局部命题}，这一点在解读检验结果时必须交代。
\end{remark}

\subsection{序数不可识别与效用标度的锚定}

定理~\ref{thm:integrability} 说 $e$ 由 $(x,\bar u)$ 唯一确定，但这句话里藏着一个循环：
$\bar u$ 本身是效用的标度，而效用标度正是不可识别的。澄清如下。

设 $e(p,\bar u)$ 是一个解，$\varphi$ 为任意严格递增函数，则
\[
 \tilde e(p,\tilde u):=e\bigl(p,\varphi^{-1}(\tilde u)\bigr)
\]
同样是解（它对应效用标度 $\tilde u=\varphi(\bar u)$），且生成同一需求。因此
\textbf{解集是 $e$ 在效用方向上的一整条重参数化轨道}。要挑出唯一代表，必须锚定标度，
最常用的三种锚定是：
\begin{enumerate}
\item \textbf{货币锚}：规定 $e(p^0,\bar u)=\bar u$，即"在基准价格下，一单位效用值一元钱"；
\item \textbf{效用锚}：规定 $u(x^0)=1$ 对某个基准束成立；
\item \textbf{指数锚}：规定 $e(p^0,\bar u)=1$，即以基准价格的生活成本为 1（生活成本指数的常用规范）。
\end{enumerate}
三种锚定给出不同的 $e$，但\emph{由 $e$ 导出的需求、Slutsky 矩阵、以及福利变化的货币度量全部相同}（\S 6.8）。
这正是"序数"的确切含义：水平不可识别，差值与曲率可识别。

\begin{remark}[一个常见误解]
"效用不可观测，所以效用函数没有意义"——这句话混淆了两个层次。
效用\emph{水平}确实不可观测也不可识别；但效用\emph{函数的曲率}（通过 $S$）完全由需求数据决定。
正因为曲率可识别，才能谈"补偿需求向下倾斜"这类可检验命题，
也才能用 $e$ 计算有经济含义的货币度量。序数性削弱的是标度，不是内容。
\end{remark}

\subsection{算例：把定理完整走一遍}

定理~\ref{thm:integrability} 的证明是纲要式的，容易让人停留在"存在性"的层面。
下面用 CD 需求把一个完整的反解过程走一遍：给定 $x$，解 Pfaff 系统，定出 $e$，
再验证 $D_p^2e$ 恰为 \S 6.5 将要算出的 Slutsky 矩阵。这条链路打通了，
可积性就不再是一个抽象定理，而是一个可执行的计算程序。

\begin{example}[从 CD 需求积分出支出函数]\label{ex:cd-invert}
设 $x_i(p,w)=\alpha_iw/p_i$，$\alpha_i>0$、$\sum_i\alpha_i=1$。Pfaff 系统 \eqref{eq:pfaff} 为
\[
 \frac{\partial e}{\partial p_i}=\frac{\alpha_i\,e}{p_i},\qquad i=1,\ldots,n .
\]
第 $i$ 个方程只含 $p_i$ 与 $e$，可分离变量：
\[
 \frac{\mathrm de}{e}=\alpha_i\,\frac{\mathrm dp_i}{p_i}
 \;\Longrightarrow\;
 e=\psi_i\bigl(p_{-i},\bar u\bigr)\,p_i^{\alpha_i},
\]
其中 $\psi_i$ 不依赖 $p_i$。对 $i=1$ 得 $e=\psi_1(p_{-1})p_1^{\alpha_1}$；
代入 $i=2$ 的方程，$\psi_1$ 中与 $p_2$ 有关的部分必须给出 $\alpha_2\ln p_2$，如此逐项进行，
最后得到
\[
 e(p,\bar u)=\varphi(\bar u)\prod_{i=1}^{n}p_i^{\alpha_i}.
\]
\textbf{齐次性自动满足}：右端对 $p$ 是 $\sum_i\alpha_i=1$ 次齐次。
这一步顺带说明 $\sum_i\alpha_i=1$ 不是参数化的人为约定：
若 $\sum_i\alpha_i\ne1$，上式对 $p$ 就不是一次齐次，因而不可能是支出函数，
需求 $x_i=\alpha_iw/p_i$ 也不是齐次性所允许的需求——两条路给出同一个结论。

\textbf{定锚}：$\varphi$ 由初值决定。取基准价格 $p^0=(\alpha_1,\ldots,\alpha_n)$ 并要求
$e(p^0,\bar u)=\bar u$，得 $\varphi(\bar u)=\bar u\big/\prod_i\alpha_i^{\alpha_i}$，故
\[
 \boxed{\;e(p,\bar u)=\bar u\prod_{i=1}^{n}\Bigl(\frac{p_i}{\alpha_i}\Bigr)^{\alpha_i}\;}
\]
与 \S 6.5 由效用最大化直接推出的结果一致。（在 $p=\alpha$ 处，消费者对每种商品各消费 $\bar u$ 单位，
总支出恰好为 $\bar u$——这就是"货币锚"在这个例子里的具体含义。）

\textbf{回代验证}：由 Shephard 引理，
\[
 \frac{\partial e}{\partial p_1}=\frac{\alpha_1}{p_1}\,e
 =\frac{\alpha_1}{p_1}\cdot w=h_1(p,\bar u),
\]
与 $x_1(p,w)=\alpha_1w/p_1$ 在对偶点 $w=e(p,\bar u)$ 处吻合，对偶恒等式得到验证。
再求二阶导：
\[
 \frac{\partial^2e}{\partial p_i\partial p_j}
 =\frac{\alpha_i\alpha_j\,e}{p_ip_j}-\delta_{ij}\frac{\alpha_i\,e}{p_i^{2}},
\]
正是 \S 6.5 命题~\ref{prop:S-cd} 给出的 $S$（在 $e=w$ 处取值）。三个环节——积分、定锚、回代——
首尾闭合，说明定理的构造不是形式上的存在性论证，而是真的可以算出东西。
\end{example}

\begin{remark}[这个算例暴露的两件事]
第一，\textbf{可积性是一个算法}：解一阶偏微分方程组 $\to$ 用一个初值定常数 $\to$ 用 Shephard 引理回代检验。
对 CD 这类可分形式，逐次积分即可；对一般的 $x$，需要沿路径数值积分（这正是 \S 6.4.3 中"路径无关"的技术内容）。
第二，\textbf{锚点不是任意的}：换成别的 $p^0$ 会得到不同形式的 $e$（相差一个只依赖 $\bar u$ 的因子），
但 $S=D_p^2e$、需求、以及福利变化的货币值都不变——这与 \S 6.4.5 关于序数标度的讨论完全一致。
\end{remark}

\section{显式恢复：三个可算系统与计量约束}
% ============================================================
% 第六讲 6.5：CD、CES、LES 的显式 Slutsky 矩阵；AIDS/QUAIDS 的约束检验
% ============================================================

\subsection{Cobb--Douglas：最简单的可积系统}

取 $u(x)=\prod_{i=1}^{n}x_i^{\alpha_i}$，$\sum_i\alpha_i=1$，$\alpha_i>0$。第 5 章已给出
\[
 x_i(p,w)=\frac{\alpha_iw}{p_i},\qquad
 e(p,\bar u)=\bar u\prod_{i=1}^{n}\Bigl(\frac{p_i}{\alpha_i}\Bigr)^{\alpha_i}.
\]
代入定义~\ref{def:slutsky}（或用 $S=D_p^2e$ 直接算）：

\begin{proposition}[CD 的 Slutsky 矩阵]\label{prop:S-cd}
对 CD 需求，
\[
 \boxed{\;S_{ij}(p,w)=\frac{\alpha_i\alpha_j\,w}{p_ip_j}-\delta_{ij}\frac{\alpha_iw}{p_i^{2}}\;}
 \qquad(\delta_{ij}=1\ \text{若}\ i=j,\ \text{否则}\ 0).
\]
\end{proposition}

\begin{proof}
由 $x_i=\alpha_iw/p_i$：$\partial x_i/\partial p_j=-\delta_{ij}\alpha_iw/p_i^2$，$\partial x_i/\partial w=\alpha_i/p_i$。
代入 $S_{ij}=\partial x_i/\partial p_j+x_j\partial x_i/\partial w$ 即得。
\end{proof}

这个矩阵的每条性质都可以直接读出：对称显然；对角元
$S_{ii}=-\alpha_iw(1-\alpha_i)/p_i^2<0$（$n\ge2$）；非对角元 $\alpha_i\alpha_jw/(p_ip_j)>0$，
故\textbf{所有商品互为总替代品}；又因 $\sum_j\alpha_j=1$，行和为
$\alpha_iw/p_i\cdot\bigl(\sum_j\alpha_j/p_j\bigr)\cdot p_j-\alpha_iw/p_i^2\cdot p_i=0$。

\begin{example}[数值验证]\label{ex:cd-num}
取 $\alpha=(0.3,0.5,0.2)$、$p=(2,3,5)$、$w=60$，则 $x=(9,\,10,\,2.4)$，
\[
 S=\begin{pmatrix}
 -3.15 & \phantom{-}1.50 & 0.36\\
 \phantom{-}1.50 & -1.6667 & 0.40\\
 \phantom{-}0.36 & \phantom{-}0.40 & -0.384
 \end{pmatrix},
\]
对称，特征值为 $(-4.0842,\,-1.1165,\,0)$，全部非正，$\operatorname{rank}S=2=n-1$。
数值上与命题~\ref{prop:S-cd} 的闭式解误差在 $10^{-8}$ 量级（浮点舍入）。
\end{example}

\begin{remark}[CD 为什么"太顺"以至于不可信]
CD 的三个特征同时成立：支出份额恒定（$s_i=\alpha_i$）、所有价格弹性为 $-1$、所有财富弹性为 $1$、
所有商品互为总替代。这四条在数据中几乎从不成立：份额随收入变化（Engel 曲线非直线）、
必需品与奢侈品的弹性不同、部分商品互为互补。
因此 CD 在理论上是最方便的基准，在实证中却几乎总被拒绝——它的价值在于提供一个\emph{可解的参照系}，
而不是一个可用的描述。
\end{remark}

\subsection{CES：把替代弹性作为参数}

取 $u(x)=\bigl(\sum_i\alpha_ix_i^{\rho}\bigr)^{1/\rho}$，$\rho<1$，$\rho\ne0$。
记替代弹性 $\sigma=1/(1-\rho)>0$。需求与支出函数为
\[
 x_i(p,w)=\frac{w\,(\alpha_i/p_i)^{\sigma}}{\sum_j p_j(\alpha_j/p_j)^{\sigma}},
 \qquad
 e(p,\bar u)=\bar u\Bigl[\sum_j\alpha_j^{\sigma}p_j^{1-\sigma}\Bigr]^{1/(1-\sigma)}.
\]
Slutsky 矩阵没有 CD 那样简洁的闭式，但可以完全显式地算出并作数值检验。
表~\ref{tab:ces} 给出四组 $\sigma$ 下的结果（$p=(2,3,5)$、$w=60$、$\alpha=(0.4,0.35,0.25)$）。

\begin{table}[H]
\centering
\caption{CES 系统在不同替代弹性下的数值结果}\label{tab:ces}
\footnotesize
\begin{tabular}{cccccc}
\toprule
$\sigma$ & $x(p,w)$ & $S$ 的三个特征值 & $\operatorname{rank}S$ & 最大非对称量 & $S_{11}$\\
\midrule
$0.5$ & $(8.835,\,6.748,\,4.417)$ & $(-1.818,\,-0.764,\,0)$ & 2 & $1.0\times10^{-9}$ & $-1.558$\\
$1.0$ & $(12.00,\,7.000,\,3.000)$ & $(-4.341,\,-1.225,\,0)$ & 2 & $8.9\times10^{-10}$ & $-3.600$\\
$2.0$ & $(18.00,\,6.125,\,1.125)$ & $(-9.316,\,-1.124,\,0)$ & 2 & $8.9\times10^{-9}$ & $-7.200$\\
$4.0$ & $(25.35,\,2.935,\,0.099)$ & $(-11.052,\,-0.225,\,0)$ & 2 & $7.5\times10^{-10}$ & $-7.859$\\
\bottomrule
\end{tabular}
\end{table}
\figexplan{四组替代弹性下的 CES 需求与 Slutsky 矩阵。所有情形下 $S$ 对称（最大非对称量在浮点误差量级）、负半定（特征值非正）、秩为 $n-1=2$，与定理~\ref{thm:slutsky-necessity} 一致。$\sigma$ 增大时 $S_{11}$ 的绝对值上升：商品间越易替代，补偿需求对自身价格越敏感。$\sigma\to\infty$ 时趋于完全替代（$S_{11}\to-\infty$ 的方向），$\sigma\to0$ 时趋于 Leontief（$S_{11}\to0$）。注意 $\sigma=1$ 一行恰为 CD，此时 $S$ 与例~\ref{ex:cd-num} 结构相同（份额不同故数值不同）。}

\begin{remark}[$\sigma$ 的经济含义与它的可检验性]
$\sigma$ 是"相对价格变动 1\% 引起相对数量变动百分之几"的度量，即 Hicks 替代弹性的常数版本。
在 CES 中它是一个标量参数，因此"检验可积性"退化为"检验 $\sigma>0$"（$\rho<1$）——
只要估计出的 $\sigma$ 为正，系统自动满足对称与负半定。
这说明\textbf{函数形式的选择本身就在施加理论约束}：选 CES 等于事先承诺了可积性，
代价是放弃了检验理论的机会。这正是 \S 6.6 转向非参数方法的动机之一。
\end{remark}

\subsection{线性支出系统：与 CD 的同构}

线性支出系统（linear expenditure system, LES，Stone 1954）设定
\begin{equation}\label{eq:les}
 x_i(p,w)=b_i+\frac{a_i}{p_i}\Bigl(w-\sum_{k=1}^{n}p_kb_k\Bigr),\qquad
 \sum_i a_i=1,\quad a_i,b_i>0.
\end{equation}
其中 $b_i$ 为"生存必需量"，$A(p,w):=w-\sum_k p_kb_k$ 称为\textbf{超numerary 支出}
（supernumerary expenditure），即扣除必需支出后可用于"自由分配"的部分；
$a_i$ 是这部分支出的边际份额。LES 由 Stone 从效用函数
$u=\sum_i a_i\ln(x_i-b_i)$ 推出，是 Gorman 极化形式的最早实例（\S 6.7）。

\begin{proposition}[LES 的 Slutsky 矩阵与 CD 同构]\label{prop:S-les}
在 \eqref{eq:les} 下，
\[
 \boxed{\;S_{ij}(p,w)=A(p,w)\left[\frac{a_ia_j}{p_ip_j}-\delta_{ij}\frac{a_i}{p_i^{2}}\right]
 =A(p,w)\,S^{\mathrm{CD}}_{ij}(a,\,1)\;}
\]
其中 $S^{\mathrm{CD}}(a,1)$ 是份额为 $a$、财富为 $1$ 的 CD 系统的 Slutsky 矩阵（命题~\ref{prop:S-cd}）。
\end{proposition}

\begin{proof}
记 $A=w-\sum_kp_kb_k$。由 \eqref{eq:les}：
\[
 \frac{\partial x_i}{\partial p_j}=-\delta_{ij}\frac{a_iA}{p_i^{2}}-\frac{a_ib_j}{p_i},
 \qquad
 \frac{\partial x_i}{\partial w}=\frac{a_i}{p_i}.
\]
代入 $S_{ij}=\partial x_i/\partial p_j+x_j\,\partial x_i/\partial w$：
\begin{align*}
S_{ij}&=-\delta_{ij}\frac{a_iA}{p_i^{2}}-\frac{a_ib_j}{p_i}
 +\frac{a_i}{p_i}\Bigl(b_j+\frac{a_jA}{p_j}\Bigr)\\
&=-\delta_{ij}\frac{a_iA}{p_i^{2}}+\frac{a_ia_jA}{p_ip_j}
=A\left[\frac{a_ia_j}{p_ip_j}-\delta_{ij}\frac{a_i}{p_i^{2}}\right].
\end{align*}
方括号内正是命题~\ref{prop:S-cd} 在 $\alpha=a$、$w=1$ 处的表达式。
\end{proof}

\begin{example}[数值核对]\label{ex:les-num}
取 $a=(0.5,0.3,0.2)$、$b=(1,2,0.5)$、$p=(2,3,5)$、$w=40$。
则 $\sum_kp_kb_k=10.5$，$A=29.5$，$x=(8.375,\,4.95,\,1.68)$，
\[
 S=\begin{pmatrix}
 -1.8438 & 0.7375 & 0.2950\\
 \phantom{-}0.7375 & -0.6883 & 0.1180\\
 \phantom{-}0.2950 & 0.1180 & -0.1888
 \end{pmatrix}
 =29.5\times S^{\mathrm{CD}}(a,1).
\]
直接按定义作数值微分与命题~\ref{prop:S-les} 的闭式解相差 $6\times10^{-9}$，
确认同构关系成立。
\end{example}

\begin{remark}[同构的经济含义与一个必须注意的边界]
命题~\ref{prop:S-les} 说的是：\textbf{LES 的替代效应与 CD 完全一样，区别只在标量 $A$}。
经济上这很自然——LES 把消费者分成两个阶段：先用 $w$ 购买必需量 $b$（"生存"），
再把剩余 $A$ 按 CD 规则分配（"享受"）。因此价格变化的替代效应只作用于第二阶段，
其结构与 CD 相同，幅度被 $A$ 缩放。

由此立刻得到一个可检验的边界条件：$S$ 负半定当且仅当 $A\ge0$（因为 $S^{\mathrm{CD}}$ 本身负半定）。
\textbf{若 $A<0$，即财富不足以购买必需量，LES 失去可积性}：
此时 \eqref{eq:les} 给出的 $x$ 仍满足会计恒等式，却不可能来自任何偏好。
实证中若样本包含贫困家庭而 $b$ 被估计得过高，就会出现 $A<0$，模型在福利评价上随之失效。
这是"参数化越强、越容易在数据边缘崩溃"的典型例子。
\end{remark}

三个系统并排比较，便于看清各自的约束强度：

\begin{center}
\footnotesize
\begin{tabular}{l>{\raggedright\arraybackslash}p{4.2cm}>{\raggedright\arraybackslash}p{4.5cm}>{\raggedright\arraybackslash}p{3.9cm}}
\toprule
系统 & \multicolumn{1}{c}{$x_i(p,w)$} & \multicolumn{1}{c}{$S$ 的结构} & \multicolumn{1}{c}{隐含的行为限制}\\
\midrule
CD & $\alpha_iw/p_i$ & $a_ia_jw/(p_ip_j)-\delta_{ij}a_iw/p_i^2$ & 份额恒定、弹性全为 1\\
CES & $w(a_i/p_i)^{\sigma}/\!\sum_jp_j(a_j/p_j)^{\sigma}$ & 无简式，恒对称负半定 & 替代弹性恒为 $\sigma$\\
LES & $b_i+(a_i/p_i)A$ & $A\cdot S^{\mathrm{CD}}(a,1)$ & 线性 Engel 曲线、分两阶段\\
\bottomrule
\end{tabular}
\end{center}
\figexplan{三个系统的对照。三者的约束强度依次为：CD 最强（份额恒定，等于把全部弹性钉死）、CES 居中（只钉死替代弹性）、LES 最弱（只要求 Engel 曲线线性）。注意三者的 $S$ 都自动对称——这不是巧合，而是因为它们都由某个具体的效用函数推出；反过来说，\textbf{任何参数化系统只要写成显式效用函数再求导，可积性就自动满足，理论检验的机会随之消失}。}

\subsection{从理论约束到计量约束：AIDS 的检验}

实证中最常用的不是上面三个系统，而是 Deaton--Muellbauer（1980）的几乎理想需求系统（AIDS）：
\[
 s_i=\alpha_i+\sum_{j=1}^{n}\gamma_{ij}\ln p_j+\beta_i\ln\!\Bigl(\frac{w}{P(p)}\Bigr),
 \qquad
 \ln P(p)=\alpha_0+\sum_k\alpha_k\ln p_k+\tfrac12\sum_{k}\sum_{l}\gamma_{kl}\ln p_k\ln p_l .
\]
它是"任意"需求系统（能二阶逼近任意需求函数），但理论约束以线性参数约束的形式出现：

\begin{center}
\begin{tabular}{lll}
\toprule
理论要求 & AIDS 参数约束 & 检验方式\\
\midrule
加总（adding-up） & $\sum_i\alpha_i=1,\ \sum_i\gamma_{ij}=0,\ \sum_i\beta_i=0$ & 估计时直接施加\\
零次齐次 & $\sum_j\gamma_{ij}=0,\ \forall i$ & Wald 检验\\
\textbf{Slutsky 对称} & $\gamma_{ij}=\gamma_{ji},\ \forall i\ne j$ & Wald / LR 检验\\
负半定 & 由 $\gamma$ 组装的 Slutsky 矩阵负半定 & 特征值检验\\
\bottomrule
\end{tabular}
\end{center}
\figexplan{AIDS 的四类约束及其检验方式。前两行是会计恒等式，通常在估计中直接施加（否则参数不可解释）；第三行是对称性，是真正的理论检验对象，自由度为 $n(n-1)/2$；第四行是曲率，检验起来最麻烦（涉及矩阵的惯性），文献中常以"施加约束后估计"代替"检验"。Banks--Blundell--Lewbel（1997）的 QUAIDS 在 $\ln(w/P)$ 上加二次项，约束形式不变。}

\begin{remark}[实证中对称性常被拒绝，怎么理解]
跨国的需求系统研究里，对称性检验的 $p$ 值经常落在拒绝域。可能的原因有四：
(1) 加总偏误——用总量数据估计个体偏好，而加总本身不保 Slutsky 条件（\S 6.7）；
(2) 动态调整——当期需求依赖上期消费（习惯形成），静态模型被误设；
(3) 价格测量误差——价格指数构造带来的衰减偏误；
(4) 偏好异质——不同家庭的参数不同，混合样本的"平均需求"不再是任何个体的需求。
这四条都不意味着"消费者不理性"，而意味着\textbf{模型的可积性条件与数据的生成过程不匹配}。
把拒绝直接读成"理论被推翻"是过度解读；反过来，把约束强行施加后宣称"模型通过检验"同样是循环论证。
正确的做法是报告施加前与施加后的估计，并说明差异的来源。
\end{remark}

\section{非参数路径：显示偏好与 Afriat 定理}
% ============================================================
% 第六讲 6.6：WARP/SARP/GARP、三个公理的关系、Afriat 定理与效率指数
% ============================================================

\subsection{数据与两个基本关系}

上一节假定知道整个需求函数。实证中更常见的情形是：只有有限个价格—数量对。

\begin{definition}[观测数据]\label{def:data}
一组观测为 $\mathcal D=\{(p^t,x^t)\}_{t=1}^{T}$，其中 $p^t\in\R^n_{++}$、$x^t\in\R^n_+\setminus\{0\}$。
记第 $t$ 期的支出 $r^t:=p^t\cdot x^t$。
\end{definition}

注意这里\emph{没有}假定预算被花光：$r^t$ 只是实际支出，可能小于财富 $w^t$。
这使非参数方法能容纳储蓄与角点，代价是"财富"从数据中消失，只剩支出。

\begin{definition}[直接显示偏好]\label{def:direct-rp}
称 $x^s$ 在 $t$ 期\textbf{直接显示为可负担}（记 $x^t\ R^0\ x^s$），若
\[
 p^t\cdot x^s\le p^t\cdot x^t=r^t .
\]
若不等号严格（$p^t\cdot x^s<r^t$），称 $x^t$ \textbf{直接严格显示偏好于} $x^s$，记 $x^t\ P^0\ x^s$。
\end{definition}

也就是说：$x^s$ 花得起而消费者却选了 $x^t$，那么 $x^t$ 至少和 $x^s$ 一样好。
这是显示偏好纲领的全部内容——它把"偏好"从不可观测的心理对象，替换成可观测的比较关系。

把 $R^0$ 取传递闭包，得到\emph{显示偏好关系} $R$（含若干步的链）；链上只要有一步严格，
就记作 $P$。三个公理分别对 $R^0$ 与 $R$ 施加不同强度的约束。

\subsection{三个公理}

\begin{definition}[WARP、SARP、GARP]\label{def:axioms}
\begin{enumerate}
\item[\textbf{WARP}]（弱公理）若 $x^t\ R^0\ x^s$ 且 $x^t\ne x^s$，则 \emph{不可能} $x^s\ R^0\ x^t$。
\item[\textbf{SARP}]（强公理）若 $x^t\ R\ x^s$ 且 $x^t\ne x^s$，则 \emph{不可能} $x^s\ R^0\ x^t$。
\item[\textbf{GARP}]（广义公理）若 $x^t\ R\ x^s$，则 \emph{不可能} $x^s\ P^0\ x^t$。
\end{enumerate}
\end{definition}

三者的差别全在"用哪个关系做前提、用哪个关系做结论"：
WARP 前提与结论都用 $R^0$（只看两期成对比较）；
SARP 前提用传递闭包 $R$、结论仍用 $R^0$；
GARP 前提用 $R$、结论用严格关系 $P^0$——比 SARP 弱，因为结论只禁止\emph{严格}的逆偏好。

\subsection{三个公理之间的关系}

\begin{proposition}[SARP $\Rightarrow$ GARP]\label{prop:sarp-garp}
若数据满足 SARP，则必满足 GARP。
\end{proposition}

\begin{proof}
反设 GARP 失败：存在 $t,s$ 使 $x^t\ R\ x^s$ 且 $x^s\ P^0\ x^t$。
由 $P^0\subseteq R^0\subseteq R$，有 $x^s\ R\ x^t$。
于是对 $t\ne s$ 同时有 $x^t\ R\ x^s$ 与 $x^s\ R^0\ x^t$（后者由 $P^0$ 蕴含），
这正是 SARP 所禁止的。故 SARP 成立时 GARP 必成立。
\end{proof}

反向不成立，且 GARP 与 WARP 之间\emph{互不包含}。两个方向各有一个反例。

\begin{example}[GARP 成立但 WARP 失败]\label{ex:garp-not-warp}
两期、两商品，价格相同：$p^1=p^2=(1,1)$，$x^1=(1,0)$，$x^2=(0,1)$。
则 $r^1=r^2=1$，且
\[
 p^1\cdot x^2=1\le r^1,\qquad p^2\cdot x^1=1\le r^2,
\]
故 $x^1\ R^0\ x^2$ 与 $x^2\ R^0\ x^1$ 同时成立而 $x^1\ne x^2$——\textbf{WARP 失败}。
但 $P^0$ 在两个方向上都不成立（两个束的支出都恰好等于 1，没有严格不等式），
GARP 的结论要求 $x^s\ P^0\ x^t$ 才能构成违背，故 \textbf{GARP 成立}。
\end{example}

\begin{example}[WARP 成立但 GARP 失败]\label{ex:warp-not-garp}
三期、三商品：
\[
\begin{array}{c|ccc|c}
t & p^t & x^t & r^t\\
\hline
1 & (2,1,1) & (1,1,2) & 5\\
2 & (1,1,2) & (1,2,1) & 5\\
3 & (2,3,1) & (2,1,1) & 8
\end{array}
\]
逐对检验：$p^1\cdot x^2=2+1+2=5\le r^1$（可负担），而 $p^2\cdot x^1=1+1+4=6>r^2$（不可负担），
故 $1\to2$ 是单向的；同理 $p^2\cdot x^3=2+1+2=5\le r^2$ 而 $p^3\cdot x^2=2+6+1=9>r^3$，
$2\to3$ 也是单向的；$p^3\cdot x^1=2+3+2=7<r^3$，故 $x^3\ P^0\ x^1$，而
$p^1\cdot x^3=4+1+1=6>r^1$，$1\to3$ 反向不成立。
三对观测逐一检查都不违反 \textbf{WARP}。

但由 $1\to2\to3$ 的传递链，$x^1\ R\ x^3$；又 $x^3\ P^0\ x^1$。
这正是 GARP 所禁止的"循环严格偏好"，故 \textbf{GARP 失败}。
又因 SARP $\Rightarrow$ GARP（命题~\ref{prop:sarp-garp}），该数据\textbf{同时违反 SARP}：
三个束互不相同，$x^1\ R\ x^2\ R\ x^3$ 且 $x^3\ R^0\ x^1$，构成一个完整的循环。
于是这一个例子同时充当"WARP $\not\Rightarrow$ SARP"与"WARP $\not\Rightarrow$ GARP"的反例。
\end{example}

\begin{remark}[两商品是例外]
\S 6.3 已看到一个现象：维数越低，理论约束越容易被自动满足。显示偏好这边有完全平行的事实：
\textbf{在两商品情形下，WARP 与 SARP 等价}（Rose, 1958, \emph{Review of Economic Studies} 25(2), 124--125），
从而也与 GARP 等价。原因不难直观理解：两商品时预算线把正象限切成两块，
"可负担"关系受到严格的几何约束，成对比较已经足以排除任意长的循环。
因此例~\ref{ex:warp-not-garp} 必须用三种商品才成立。
这一事实有一个直接的实证含义：\textbf{若研究只涉及两类商品（如"食品"与"其他"），
用 WARP 还是 GARP 做检验没有区别}；理论的可检验内容被维数压缩了。
\end{remark}

\begin{remark}[一个实操细节：同一束重复出现]
三个公理的定义都要求 $x^t\ne x^s$，这条看似技术性的限定在数据清洗阶段很重要：
同一消费束在不同价格下重复出现时，$R^0$ 在两个方向上都成立，但三者的判定会分岔——
WARP 与 GARP 都因 $x^t=x^s$ 而豁免（不构成违背），而若把"束相同"误当作"偏好严格循环"则会误报违规。
反过来，若两期观测到同一束而价格不同，这本身就是对"需求函数"概念的一个提醒：
同一个人在两组不同价格下选了同一束，说明预算约束并未绑定，或偏好在两期间发生了变化。
合并这类重复观测应当在检验之前完成，否则检验结果取决于公理的措辞而非数据。
\end{remark}

\subsection{Afriat 定理：GARP 是可理性化的充要条件}

\begin{theorem}[Afriat, 1967；Varian, 1982]\label{thm:afriat}
对有限数据 $\mathcal D=\{(p^t,x^t)\}_{t=1}^{T}$，下列命题等价：
\begin{enumerate}
\item $\mathcal D$ 满足 GARP；
\item 存在一个连续、局部非饱和的效用函数 $u:\R^n_+\to\R$，使得对一切 $t$，
      $x^t$ 是 $u$ 在预算集 $\{x:p^t\cdot x\le p^t\cdot x^t\}$ 上的最优解；
\item 存在实数 $\{u^t\}_{t=1}^{T}$ 与 $\{\lambda^t\}_{t=1}^{T}$（$\lambda^t>0$），使 Afriat 不等式成立：
      \[
       u^s\le u^t+\lambda^t\,p^t\cdot(x^s-x^t),\qquad \forall s,t=1,\ldots,T .
      \]
\end{enumerate}
\end{theorem}

\begin{remark}[定理的三层含义]
第一，\textbf{GARP 是充要的}：它不仅是可理性化的必要条件（由显示偏好的定义直接可得），
还是充分条件。这一点远非显然——GARP 只约束成对与链式比较，却足以保证存在一个\emph{全局}的效用函数。
第二，\textbf{不需要任何函数形式}：结论是"存在某个 $u$"，而不是"存在某个 CD 或 CES"。
第三，\textbf{可计算}：Afriat 不等式是一组线性不等式，可行性可用线性规划判定，
因此检验的复杂度是多项式时间的，$T$ 上千也能算。
\end{remark}

\begin{remark}[Afriat 不等式的读法]
$\lambda^t$ 可解释为第 $t$ 期"一元钱的边际效用"：把 $u^s$ 与 $u^t$ 的差用支出差
$p^t\cdot(x^s-x^t)$ 线性化，$\lambda^t$ 就是斜率。若 $x^s$ 在 $t$ 期可负担
（$p^t\cdot(x^s-x^t)\le0$），不等式给出 $u^s\le u^t$——正是显示偏好的要求。
不等式组存在解，等价于存在一条与全部观测相容的"分片线性"效用函数；
Afriat 的构造正是用这些 $(\lambda^t,u^t)$ 拼出一个分段线性的 $u$。
\end{remark}

\begin{definition}[Afriat 效率指数]
若 GARP 被违反，可定义效率指数
\[
 e^*:=\max\bigl\{\,e\in(0,1]:\ \text{把每个}\ x^t\ \text{按比例缩减为}\ e\,x^t\ \text{后数据满足 GARP}\,\bigr\}.
\]
$e^*=1$ 当且仅当原数据满足 GARP；$e^*$ 越接近 1，数据"越接近"可理性化。
\end{definition}

\begin{remark}[指数不能当拟合优度用]
$e^*$ 常被误读为"理性程度"。它度量的是\emph{为恢复可理性化所需的最小统一缩减比例}，
是一个几何量，没有行为含义：它同时吸收了测量误差、偏好异质与真实的非理性，
三者无法从数据中分离。把它报告为"消费者有 $e^*$ 比例的理性"是范畴错误。
合理的用法只有两个：作为设定诊断（$e^*$ 明显低于 1 说明静态理性模型不适合这批数据），
以及作为稳健性比较（不同子样本、不同商品划分下的 $e^*$ 相对大小）。
\end{remark}

\subsection{两条路径的对照}

\begin{center}
\begin{tabular}{lll}
\toprule
维度 & 微分路径（\S 6.2--6.5） & 非参数路径（\S 6.6）\\
\midrule
输入 & 需求函数的解析式 & 有限观测点\\
核心条件 & $S$ 对称 $+$ 负半定 & GARP\\
输出 & 闭式的 $e(p,\bar u)$、$h$ & 一个存在性结论，无闭式\\
函数形式 & 需事先指定（CD/CES/LES/AIDS） & 不需要\\
可检验内容 & 对称性、曲率（参数约束） & 整个显示偏好结构\\
能否算福利 & 能（$e$ 已知） & 不能直接算\\
\bottomrule
\end{tabular}
\end{center}
\figexplan{两条恢复路径的对照。二者并不冲突：若数据来自一个满足可积性条件的光滑需求系统，则它必满足 GARP（取足够密的观测即可验证）；反过来，GARP 成立只保证存在\emph{某个}效用函数，不保证它光滑、可微、能生成可积的需求函数。因此非参数检验更稳健（假设更少），参数方法更有力（结论更强）。实证中的常见做法是先用非参数方法做设定筛查，再选参数系统做估计与福利计算。}

\begin{figure}[H]
\centering
\begin{tikzpicture}[scale=0.95]
  \draw[axisline] (0,0)--(6.4,0) node[below]{$x_1$};
  \draw[axisline] (0,0)--(0,5.0) node[left]{$x_2$};
  % 第 1 期预算集
  \fill[black!8] (0,0)--(4.6,0)--(0,4.6)--cycle;
  % 预算线 1：x1+x2=4.6
  \draw[curveA] (0,4.6)--(4.6,0);
  % 预算线 2：0.5x1+x2=2.2
  \draw[curveB] (0,2.2)--(4.4,0);
  % 两个观测束
  \node[nodewhite] at (2.4,2.2) {};
  \node[figlabel,above right] at (2.48,2.28) {$x^1$};
  \node[nodewhite] at (1.0,1.7) {};
  \node[figlabel,below left] at (0.92,1.60) {$x^2$};
  % 显示偏好标注（置于预算集外的空白区）
  \node[figlabel] at (3.45,3.95) {$x^1\ P^0\ x^2$};
  % 图例
  \draw[curveA] (0.1,-0.95)--(0.85,-0.95);
  \node[right,font=\scriptsize] at (0.95,-0.95)
      {第 1 期预算线 $p^1\cdot x=r^1$（浅灰为预算集）};
  \draw[curveB] (0.1,-1.42)--(0.85,-1.42);
  \node[right,font=\scriptsize] at (0.95,-1.42)
      {第 2 期预算线 $p^2\cdot x=r^2$};
\end{tikzpicture}
\caption{直接显示偏好的几何}
\figexplan{实线为第 1 期预算线，浅灰区域为其预算集；虚线为第 2 期预算线。$x^2$ 落在第 1 期预算集内部（$p^1\cdot x^2<r^1$），而消费者在第 1 期选了 $x^1$，故 $x^1$ 严格显示偏好于 $x^2$。反过来 $x^1$ 落在第 2 期预算集之外（$p^2\cdot x^1>r^2$），故反向的显示偏好不成立——这对观测不违反 WARP。所有公理的内容都可以在这样的图中读出：违背发生在"两条预算线互相穿过对方所选的束"之时。}
\label{fig:garp}
\end{figure}

\section{加总问题：代表性消费者什么时候存在}
% ============================================================
% 第六讲 6.7：加总的负结果、加总 Slutsky 的分解、Gorman 极化形式
% ============================================================

\subsection{问题的实质}

至此讨论的都是单个消费者。实证研究几乎总是用总量数据：把全国的消费支出当"一个消费者"的支出，
把总量财富当"他的"财富，然后估计一个需求系统、算一个福利变化。这一步跳跃需要理由。

设经济中有 $I$ 个消费者，第 $i$ 个的需求为 $x^i(p,w_i)$，总量需求为
\[
 X(p,w_1,\ldots,w_I)=\sum_{i=1}^{I}x^i(p,w_i).
\]
问题有两层，必须分开问：
\begin{enumerate}
\item \textbf{统计层面}：$X$ 能否写成 $X(p,W)$ 的形式，其中 $W=\sum_iw_i$？
      即总量需求是否只依赖总财富、不依赖其分配？
\item \textbf{规范层面}：若第一层成立，能否进一步把 $X$ 当作\emph{某个}偏好生成的个体需求，
      从而用 $e$ 计算"社会"的福利变化？
\end{enumerate}
第一层不成立时，代表性消费者连定义都没有；第一层成立而第二层不成立时，福利评价仍然没有依据。
两层的条件都不平凡。

\subsection{一个负结果：即使偏好相同也很难加总}

先看最有利的情形——所有消费者偏好\emph{相同}。此时加总似乎是自动的：同一函数，求和即可。
但"偏好相同"并不意味着"需求相同"，因为财富不同，而需求对财富的依赖方式可能不同。

\begin{example}[同偏好、不同财富，仍不能加总]\label{ex:agg-fail}
设两人都有 CD 偏好但份额不同（这是"偏好相同"的最宽松读法；若份额也相同则见下文注）。
取 $\alpha_1=0.7$、$\alpha_2=0.3$，$n=2$，$p=(1,1)$，总财富 $W=10$ 固定。
两人的商品 1 需求为 $x^1_1=0.7w_1/p_1$、$x^2_1=0.3w_2/p_1$，故
\[
 X_1=\frac{0.7w_1+0.3w_2}{p_1}.
\]
在三种分配下：
\begin{center}
\begin{tabular}{lccc}
\toprule
分配 $(w_1,w_2)$ & $(5,5)$ & $(10,0)$ & $(0,10)$\\
\midrule
$X_1$（$p_1=1$） & $5.0$ & $7.0$ & $3.0$\\
\bottomrule
\end{tabular}
\end{center}
总财富、价格、偏好全部不变，总量需求却在 $[3,7]$ 之间大幅移动。
故 $X$ \textbf{不是} $(p,W)$ 的函数，代表性消费者不存在。
\end{example}

\begin{remark}[这个例子为什么值得强调]
CD 是"最容易加总"的偏好：需求对财富线性、对价格弹性恒为 1、支出份额恒定。
如果连 CD 在份额不同时都无法加总，一般情形只会更糟。
关键原因是：加总把 $w_i$ 消掉的同时，也把"谁有多少钱"这一信息消掉了，
而需求恰恰依赖这一信息。只要个体需求对财富的斜率不同（$\partial x^1_1/\partial w_1\ne\partial x^2_1/\partial w_2$），
总量就依赖分配。反之，若所有个体的财富斜率相同（$0.7=0.3$），加总自动成立——
这正是 Gorman 条件的雏形（\S 6.7.4）。
\end{remark}

\subsection{强行加总的代价：加总 Slutsky 矩阵多出一项}

假设研究者无视上述问题，直接用总量数据构造 Slutsky 矩阵。这样做会得到什么？
把总量需求对 $p_j$ 与 $W$ 求导，并设每个个体的财富同比例增加（$\partial w_i/\partial W=1$）：

\begin{proposition}[加总 Slutsky 矩阵的分解]\label{prop:S-agg}
设 $X(p,W)=\sum_ix^i(p,w_i)$（假设存在），则
\[
 \boxed{\;S^{\mathrm{agg}}(p,W)
 =\sum_{i=1}^{I}S^{i}(p,w_i)
 +\underbrace{\sum_{i=1}^{I}\sum_{k\ne i}\Bigl(\frac{\partial x^{i}}{\partial w_i}\Bigr)x^{k\,\top}}_{\text{分配项}}\;}
\]
其中 $S^i$ 为第 $i$ 个个体的 Slutsky 矩阵。
\end{proposition}

\begin{proof}
由 $X_j=\sum_ix^i_j$：
\[
 \frac{\partial X_j}{\partial p_l}=\sum_i\frac{\partial x^i_j}{\partial p_l},
 \qquad
 \frac{\partial X_j}{\partial W}=\sum_i\frac{\partial x^i_j}{\partial w_i}.
\]
代入 $S^{\mathrm{agg}}_{jl}=\partial X_j/\partial p_l+(\partial X_j/\partial W)X_l$：
\[
 S^{\mathrm{agg}}_{jl}=\sum_i\frac{\partial x^i_j}{\partial p_l}
 +\Bigl(\sum_i\frac{\partial x^i_j}{\partial w_i}\Bigr)\Bigl(\sum_kx^k_l\Bigr)
 =\sum_i\underbrace{\Bigl[\frac{\partial x^i_j}{\partial p_l}+\frac{\partial x^i_j}{\partial w_i}x^i_l\Bigr]}_{=S^i_{jl}}
 +\sum_i\sum_{k\ne i}\frac{\partial x^i_j}{\partial w_i}x^k_l . \qedhere
\]
\end{proof}

\begin{remark}[分配项的读法]
分配项把"$i$ 的财富效应"乘上"$k$ 的消费量"。它的存在意味着：
\textbf{总量层面的替代效应，混入了跨人的财富再分配效应}。这带来三个后果：
(1) 即使每个 $S^i$ 都对称，$S^{\mathrm{agg}}$ 一般\textbf{不对称}——分配项本身通常不对称；
(2) $S^{\mathrm{agg}}$ 的负半定性也没有保证（分配项可正可负）；
(3) 因此用总量数据做 Slutsky 对称性检验，拒绝的可能不是消费者理论，而是加总假设本身。
这与 \S 6.5 末尾提到的"加总偏误"是同一件事的代数表达。
\end{remark}

\subsection{Gorman 极化形式：精确加总的充要条件}

什么条件下加总成立？Gorman（1953）给出了答案，其形式之美值得单独一节。

\begin{definition}[Gorman 极化形式]
称一组偏好满足\textbf{极化形式}，若其支出函数可写成
\[
 e_i(p,u_i)=c_i(p)+d(p)\,u_i,\qquad i=1,\ldots,I,
\]
其中 $c_i,d$ 对 $p$ 一次齐次、凹、非减，且 $d(p)>0$ 对\emph{所有} $i$ \textbf{相同}。
\end{definition}

\begin{proposition}[精确加总]\label{prop:gorman}
若所有个体的支出函数取 Gorman 极化形式，则总量需求可由一个代表性消费者生成：
\[
 E(p,U)=\sum_{i=1}^{I}c_i(p)+d(p)\,U,\qquad U:=\sum_{i=1}^{I}u_i,
\]
且 $X(p,W)=\nabla_pE\bigl(p,U(p,W)\bigr)$，其中 $U(p,W)$ 由 $E(p,U)=W$ 反解。
\end{proposition}

\begin{proof}
由 $e_i=c_i+du_i$ 与 Shephard 引理，$x^i=\nabla_pe_i=\nabla_pc_i+u_i\nabla_pd$。求和：
\[
 X=\sum_i\nabla_pc_i+\Bigl(\sum_iu_i\Bigr)\nabla_pd
 =\nabla_p\Bigl[\sum_ic_i(p)+d(p)U\Bigr]=\nabla_pE(p,U).
\]
所有权函数对 $p$ 一次齐次、凹，故 $E$ 亦然，从而是合法的支出函数；
$E(p,U)=W$ 对 $U$ 严格递增（因 $d>0$），可反解出 $U(p,W)$。
\end{proof}

\begin{example}[LES 属于 Gorman 形式]
LES 的支出函数为
\[
 e_i(p,u_i)=p\cdot b_i+e^{u_i}\prod_{k}p_k^{a_k},
\]
取 $c_i(p)=p\cdot b_i$、$d(p)=\prod_kp_k^{a_k}$，正是极化形式——注意 $d$ 与 $i$ 无关，
因为所有消费者的边际份额 $a$ 相同，差异只体现在必需量 $b_i$ 上。
于是总量支出为
\[
 E(p,U)=\sum_ip\cdot b_i+d(p)\,U=p\cdot\Bigl(\sum_ib_i\Bigr)+d(p)\,U,
\]
即\textbf{总量也服从一个 LES}，其"必需量"是各家必需量之和。这是 LES 在实证中被长期使用的真正原因：
它是少数几个既灵活、又\emph{精确可加总}的函数形式之一。
\end{example}

\begin{remark}[Gorman 条件的实质：共同的财富斜率]
把 $e_i=c_i+du_i$ 对 $u_i$ 求导：$\partial e_i/\partial u_i=d(p)$，与 $i$ 无关。
由包络定理，$\partial e_i/\partial u_i=1/\lambda_i$（$\lambda_i$ 为 $i$ 的财富边际效用），
故 Gorman 条件等价于说：\textbf{所有人的财富边际效用之比与财富分配无关}——
或者说，$\partial x^i/\partial w_i$ 有一个共同的（按 $d$ 缩放的）结构。
这正好解释了例~\ref{ex:agg-fail}：那里两人的财富斜率 $0.7/p_1$ 与 $0.3/p_1$ 不同，条件不满足，加总失败。
\end{remark}

\begin{remark}[Gorman 加总的代价：基数效用]
命题~\ref{prop:gorman} 中的 $U=\sum_iu_i$ 是效用的\emph{水平}之和。
这只在效用可作基数比较时才有意义——而 \S 6.1 已经说明，效用水平在个体层面不可识别。
于是出现一个悖论式的结论：\textbf{代表性消费者的存在性依赖于一个不可识别的量}。
实践中这意味着：Gorman 条件可以保证"总量需求像一个理性消费者"，但\emph{不能}保证
"这个代表性消费者的福利变化等于个体福利变化之和"。用总量数据算出的 CV/EV，
在分配发生大幅变动时可能与人头加总的福利结果符号相反。
这是 \S 6.8 之后必须交代的一条边界。
\end{remark}

把本节的结果收成一张表：

\begin{center}
\small
\begin{tabular}{lll}
\toprule
情形 & 总量需求是否只依赖 $(p,W)$ & 能否作福利评价\\
\midrule
偏好同质、财富斜率相同（Gorman 特例） & 是 & 需基数效用\\
偏好同质、财富斜率不同 & 否（例~\ref{ex:agg-fail}） & 否\\
Gorman 极化形式、$d(p)$ 相同 & 是（命题~\ref{prop:gorman}） & 需基数效用\\
一般异质偏好 & 否 & 否\\
\bottomrule
\end{tabular}
\end{center}
\figexplan{加总的三档结论。第一档是"同质偏好 $+$ 共同财富斜率"，加总自动成立；第二档是 Gorman 极化形式，加总成立但需要$\sum u_i$ 可比；第三档是一般情形，加总不成立。实证中使用总量数据时，实际站在哪一档通常无法从数据中检验——这是总量分析最根本的识别困难，比 Slutsky 检验的通过与否更值得在论文中交代。}

\section{福利评价：补偿变差、等价变差与消费者剩余}
% ============================================================
% 第六讲 6.8：CV/EV 的定义、与 Marshallian 剩余的夹逼、路径依赖
% ============================================================

\subsection{两个货币度量}

恢复出 $e$ 之后，福利评价就变成一个代入运算。考虑价格从 $p^0$ 变为 $p^1$，财富 $w$ 不变。
记 $u^0=v(p^0,w)$、$u^1=v(p^1,w)$。两个标准的货币度量是：

\begin{definition}[补偿变差与等价变差]\label{def:cv-ev}
\begin{align*}
\textbf{补偿变差（CV）}:\quad & \mathrm{CV}:=e(p^1,u^0)-e(p^1,u^1)=e(p^1,u^0)-w,\\
\textbf{等价变差（EV）}:\quad & \mathrm{EV}:=e(p^0,u^0)-e(p^0,u^1)=w-e(p^0,u^1).
\end{align*}
\end{definition}

两个定义只在"用哪一组价格做尺子"上不同：
CV 用\emph{新}价格衡量"要补多少钱才能回到原来的效用"，是事后补偿；
EV 用\emph{旧}价格衡量"多少收入变化与这次价格变化等价"，是事前等价。
对价格上涨（$u^1<u^0$），两者都为正：

\begin{proposition}[积分表示]\label{prop:cv-ev-integral}
设价格只变第 1 种商品（$p^0_1<p^1_1$，其余价格不变），则
\[
 \mathrm{CV}=\int_{p^0_1}^{p^1_1}h_1(p_1,u^0)\,\mathrm dp_1,
 \qquad
 \mathrm{EV}=\int_{p^0_1}^{p^1_1}h_1(p_1,u^1)\,\mathrm dp_1,
\]
两者均严格为正。
\end{proposition}

\begin{proof}
由 $e$ 对 $p$ 一次齐次且 $\nabla_pe=h$，沿"只改 $p_1$"的路径积分：
$e(p^1,u^0)-e(p^0,u^0)=\int_{p^0_1}^{p^1_1}h_1(p_1,u^0)\mathrm dp_1$。
而 $e(p^0,u^0)=w$，故第一式成立。第二式同理由 $e(p^0,u^1)$ 与 $e(p^1,u^1)=w$ 相减得到。
正性来自 $h_1>0$ 与积分方向。
\end{proof}

\subsection{与 Marshallian 消费者剩余的关系}

教科书中更常见的是 Marshallian 消费者剩余：
\[
 \Delta\mathrm{CS}:=-\int_{p^0_1}^{p^1_1}x_1(p_1,w)\,\mathrm dp_1
 \qquad\bigl(\text{价格上升时为负，其绝对值记}\ |\Delta\mathrm{CS}|\bigr).
\]
它与 CV、EV 的关系可以精确地"夹"出来。

\begin{proposition}[夹逼]\label{prop:sandwich}
设商品 1 为正常品，$p^0_1<p^1_1$，其余价格不变。则
\[
 \boxed{\;\mathrm{EV}\;\le\;|\Delta\mathrm{CS}|\;\le\;\mathrm{CV}.\;}
\]
\end{proposition}

\begin{proof}
对 $p_1\in[p^0_1,p^1_1]$，Marshallian 需求满足
$x_1(p_1,w)=h_1\bigl(p_1,\,v(p_1,w)\bigr)$（对偶恒等式）。
因价格上升，$v(p_1,w)$ 从 $u^0$ 单调降到 $u^1$，故 $u^1\le v(p_1,w)\le u^0$。
商品 1 正常意味着 $h_1$ 对效用水平递增，于是对区间内每个 $p_1$：
\[
 h_1(p_1,u^1)\;\le\;x_1(p_1,w)\;\le\;h_1(p_1,u^0).
\]
在 $[p^0_1,p^1_1]$ 上积分（积分区间长度为正），并用命题~\ref{prop:cv-ev-integral} 即得结论。
\end{proof}

\begin{remark}[夹逼的适用范围]
三个条件缺一不可：(i) 只变一种价格；(ii) 商品为正常品；(iii) 价格变化方向明确。
若商品为劣等品，$h_1$ 对 $u$ 递减，不等号方向反转，夹逼失效；
若同时变多种价格，$\Delta\mathrm{CS}$ 变成线积分，见 \S 6.8.4。
文献中"消费者剩余可以作为 CV/EV 的良好近似"这一说法，依赖的正是命题~\ref{prop:sandwich}——
它给出的是\emph{有界}而非"近似相等"，界的宽度取决于价格变化幅度与需求曲线的曲率。
价格变化越大，三条曲线围出的区域越宽，近似越差。
\end{remark}

\begin{example}[CD 情形：三个数都能算出来]\label{ex:cd-welfare}
取 $u=x_1^{0.3}x_2^{0.7}$，$p^0=(2,3)$、$p^1=(4,3)$、$w=60$（只变 $p_1$，涨一倍）。
由 $v(p,w)=w\prod_i(\alpha_i/p_i)^{\alpha_i}$ 与 $e(p,u)=u\prod_i(p_i/\alpha_i)^{\alpha_i}$：
\[
 u^0=12.2620,\qquad u^1=9.9599 .
\]
三个度量：
\begin{align*}
\mathrm{CV}&=w\Bigl[\frac{\prod_i(p_i^1)^{\alpha_i}}{\prod_i(p_i^0)^{\alpha_i}}-1\Bigr]
 =60\bigl(2^{0.3}-1\bigr)=13.8687,\\
\mathrm{EV}&=w\Bigl[1-\frac{\prod_i(p_i^0)^{\alpha_i}}{\prod_i(p_i^1)^{\alpha_i}}\Bigr]
 =60\bigl(1-2^{-0.3}\bigr)=11.2649,\\
|\Delta\mathrm{CS}|&=\alpha_1w\ln\frac{p^0_1}{p^1_1}=0.3\times60\times\ln2=12.4766 .
\end{align*}
于是
\[
 11.2649\;<\;12.4766\;<\;13.8687,
\]
即 $\mathrm{EV}<|\Delta\mathrm{CS}|<\mathrm{CV}$，与命题~\ref{prop:sandwich} 一致。
注意三个数的差异达到 20\% 以上（$13.87$ 对 $11.26$）：在价格翻倍的场合，
用消费者剩余替代 CV 会低估补偿需求约 10\%，这个量级在政策评估中不可忽略。
\end{example}

\begin{figure}[H]
\centering
\begin{tikzpicture}[scale=1.0]
  % 坐标轴：横轴数量，纵轴价格
  \draw[axisline] (0,0)--(7.0,0) node[below]{$x_1$};
  \draw[axisline] (0,0)--(0,5.2) node[left]{$p_1$};
  % 三个嵌套积分区域：由外到内 CV、|ΔCS|、EV
  \fill[black!8]  (0,1.2)--(0,4.4)--(1.76,4.4)--(5.28,1.2)--cycle;
  \fill[black!22] (0,1.2)--(0,4.4)--(1.60,4.4)--(4.80,1.2)--cycle;
  \fill[black!42] (0,1.2)--(0,4.4)--(1.44,4.4)--(4.32,1.2)--cycle;
  % 两条价格水平线
  \draw[dashed,thick] (0,4.4)--(6.1,4.4);
  \node[figlabel,left] at (0,4.4) {$p^1_1$};
  \draw[dashed,thick] (0,1.2)--(6.1,1.2);
  \node[figlabel,left] at (0,1.2) {$p^0_1$};
  % 三条需求曲线（只在价格区间内绘制）
  \draw[curveA,domain=1.05:4.5,smooth] plot ({6.6-1.1*\x},{\x});
  \draw[curveB,domain=1.05:4.5,smooth] plot ({6.0-1.0*\x},{\x});
  \draw[curveC,domain=1.05:4.5,smooth] plot ({5.4-0.9*\x},{\x});
  % 图例
  \draw[curveA] (0.1,-0.95)--(0.85,-0.95);
  \node[right,font=\scriptsize] at (0.95,-0.95)
      {$h_1(p_1,u^0)$：$\mathrm{CV}$ 的积分路径（最外）};
  \draw[curveB] (0.1,-1.42)--(0.85,-1.42);
  \node[right,font=\scriptsize] at (0.95,-1.42)
      {$x_1(p_1,w)$：$|\Delta\mathrm{CS}|$ 的积分路径（居中）};
  \draw[curveC] (0.1,-1.89)--(0.85,-1.89);
  \node[right,font=\scriptsize] at (0.95,-1.89)
      {$h_1(p_1,u^1)$：$\mathrm{EV}$ 的积分路径（最内）};
\end{tikzpicture}
\caption{补偿变差、等价变差与消费者剩余}
\figexplan{横轴为商品 1 的数量，纵轴为其价格。三条曲线自上而下分别是 Hicksian 需求 $h_1(p_1,u^0)$（实线）、Marshallian 需求 $x_1(p_1,w)$（虚线）、Hicksian 需求 $h_1(p_1,u^1)$（点线）——正常品前提下它们的顺序固定，这正是命题~\ref{prop:sandwich} 的几何内容。价格由 $p^0_1$ 升至 $p^1_1$，由价格轴量至三条曲线所围出的三个嵌套区域，其面积依次为 $\mathrm{CV}$（最外）、$|\Delta\mathrm{CS}|$（居中）、$\mathrm{EV}$（最内）；三者之差即图中可见的两条窄带。曲线画成直线只是为了便于标注面积，真实情形下它们向下凸，三者的差距随价格变动幅度增大。}
\label{fig:cv-ev-cs}
\end{figure}

\subsection{多价格同时变动：路径依赖}

命题~\ref{prop:sandwich} 要求"只变一种价格"。一旦多种价格同时变动，$\Delta\mathrm{CS}$ 必须写成线积分
\[
 \Delta\mathrm{CS}=-\int_{\Gamma}\sum_{i=1}^{n}x_i(p,w)\,\mathrm dp_i ,
\]
而 Marshallian 需求的交叉价格效应一般\emph{不对称}（$\partial x_i/\partial p_j\ne\partial x_j/\partial p_i$），
故这个积分的值与路径 $\Gamma$ 有关。

\begin{proposition}[CV/EV 无路径依赖]\label{prop:no-path}
$\mathrm{CV}$ 与 $\mathrm{EV}$ 与价格调整路径无关；$\Delta\mathrm{CS}$ 一般依赖路径。
\end{proposition}

\begin{proof}
$\mathrm{CV}=e(p^1,u^0)-w$ 只是两个支出函数值之差，$e$ 是良定义的函数，与路径无关。
$\mathrm{EV}$ 同理。而 $\Delta\mathrm{CS}$ 的线积分是 1-形式 $\omega=-\sum_ix_i\,\mathrm dp_i$ 的积分，
该形式闭（$\mathrm d\omega=0$）当且仅当 $\partial x_i/\partial p_j=\partial x_j/\partial p_i$，
即 Marshallian 交叉导数对称——这在一般情形下不成立（对称的是 Hicksian 一侧的 $S$）。
故 $\omega$ 非闭，积分依赖路径。
\end{proof}

\begin{remark}[这是 CV/EV 最重要的技术优势]
许多人把 CV/EV 视为"更精细"的度量，理由说成"它考虑了一般均衡"之类，这是误解。
真正的理由是：\textbf{CV/EV 建立在支出函数这个良定义的势函数上，因此对多价格变动自动路径无关；
消费者剩余建立在 Marshallian 需求这个非保守场上，只在"单一价格变动"这一特殊情形才恰好良定义。}
这也解释了为什么可积性（\S 6.4）是福利分析的前提：没有 $e$，就只有一个非保守的向量场，
任何"面积"都必须先指定路径。
\end{remark}

\subsection{加总与可加性：两条不可逾越的边界}

即使 $e$ 已恢复，福利度量还有两条边界必须交代。

\notehead{边界一：CV/EV 不可跨人相加。}

$e(p,u)$ 的数值依赖效用标度。甲的 $\mathrm{CV}=100$ 与乙的 $\mathrm{CV}=100$ 之和没有意义，
除非两人的效用被置于同一基数尺度上——而这正是 \S 6.1 观察一所否定的。
因此"总补偿变差"不是一个福利概念，只是一个会计加总。要谈社会偏好，必须引入显性的社会福利函数，
那是规范选择，不能从个体需求数据中推出。

\notehead{边界二：消费者剩余可加，但没有福利含义。}

$\Delta\mathrm{CS}$ 以货币为单位、可以跨人相加，看起来很适合做成本—收益分析。
但由命题~\ref{prop:sandwich}，它在个体层面只是 CV 与 EV 之间的一个有界近似；
在加总层面，它既不等于总 CV 也不等于总 EV，且其与"真实"福利变化的偏差随价格变动幅度与异质性放大。
实践中用 $\Delta\mathrm{CS}$ 做政策评估是可以的，前提是明确它是一个\emph{近似指标}，
并在论文中给出敏感性分析（例如同时报告 CV/EV 上界）。

\begin{remark}[一个常见的误用]
用总量数据估出一个需求系统，代入公式算出 CV，然后宣称"该政策使社会福利提高 X 亿元"。
这句话至少跨过了三道未验证的门槛：加总成立否（\S 6.7）、效用可比否（观察一）、
以及代表性消费者的 $e$ 是否等于个体 $e$ 的加总（Gorman 条件）。
规范的写法是把结论限制在"代表性消费者的补偿变差"这一明确对象上，
并说明它与社会福利的差距取决于分配效应。
\end{remark}

\section{从数据到福利：一条可执行的工作流}
% ============================================================
% 第六讲 6.9：六步工作流、三类识别困难、检验清单
% ============================================================

\subsection{六步工作流}

前面八节的技术内容可以压缩成一条流水线。下表把每一步的输入、输出与\emph{可能失败的地方}并列，
后者比前者更值得记住。

\begin{center}
\footnotesize
\begin{tabular}{@{}c>{\raggedright\arraybackslash}p{5.2cm}>{\raggedright\arraybackslash}p{3.0cm}>{\raggedright\arraybackslash}p{4.3cm}@{}}
\toprule
步 & \multicolumn{1}{c}{操作} & \multicolumn{1}{c}{输出} & \multicolumn{1}{c}{常见失败}\\
\midrule
1 & 清洗数据、合并重复观测、定义商品与价格 & $\{(p^t,x^t)\}$ & 商品划分随数据可得性而变\\
2 & 非参数筛查：检验 GARP、算 Afriat 指数 & 是否可理性化 & 忽略角点与动态调整\\
3 & 选择函数形式（CD/CES/LES/AIDS/QUAIDS） & 参数化系统 & 形式本身已施加可积性\\
4 & 估计参数（施加加总与齐次约束） & 参数估计值 & 价格内生、测量误差\\
5 & 检验对称性与曲率 & 检验统计量、$p$ 值 & 把拒绝直接读成"理论错"\\
6 & 由 $e$ 计算 CV/EV，做敏感性分析 & 福利度量及区间 & 忽略加总与效用可比\\
\bottomrule
\end{tabular}
\end{center}
\figexplan{六步工作流。第 2 步与第 3 步是互斥的两条路：若走非参数路线，第 3--5 步可以跳过，但第 6 步无法执行（没有 $e$）；若走参数路线，第 2 步的 GARP 检验仍值得做一次作为设定诊断。表中第四列列出的是每一步最容易被忽视的失败模式，它们大多不是估计技术问题，而是"数据生成过程与模型假设不匹配"的问题。}

\subsection{第 2 步：非参数筛查怎么做}

GARP 检验的计算量随 $T$ 增长但不是瓶颈，真正的难点在两处。

\notehead{其一，价格与数量的口径必须一致。}
观测必须是"同一个决策单位、同一时期、面对同一组价格"的选择。用年度均价配年度总消费，
等于假定年内价格不变；用地区均价配家庭消费，等于假定家庭面对的是地区均价。
两种做法都会人为地\emph{减少}价格变异，从而让 GARP 更容易通过——
\textbf{检验的功效被数据构造悄悄削弱了}。这一点在论文中应当明说。

\notehead{其二，角点与耐用品的处理。}
\S 6.2 的 (A4) 假定内点，而真实数据中大量商品是零消费（不买酒、不买耐用品）。
非参数方法本身不需要内点假设，但"零消费"会使 $R^0$ 的判定失去意义（$p^t\cdot x^s$ 中缺失的分量无法补）。
常见处理是加总成大类（食品、衣着、居住……）以降低零消费比例，代价是商品数 $n$ 变小，
可检验内容随之减少（\S 6.3 的表）。

\subsection{第 4--5 步：把拒绝读对}

\begin{remark}[三种拒绝，三种解释]
设对称性检验被拒绝。至少有三种互不排斥的解释，而数据本身无法区分它们：
\begin{enumerate}
\item \textbf{偏好异质}：样本混合了参数不同的家庭，混合后的"平均需求"不是任何个体的需求（\S 6.7）；
\item \textbf{动态因素}：习惯形成、耐用品存量调整，使当期需求依赖历史，静态模型被误设；
\item \textbf{测量误差}：价格指数与数量的构造误差，在对称性检验中不抵消而累积。
\end{enumerate}
因此规范的表述是"在某某设定下对称性被拒绝"，而不是"消费者理论被证伪"。
反过来，把约束强行施加后再报告"模型通过检验"则是循环论证——
检验的对象已经被假设掉了。可行的折中是：报告无约束估计、约束估计与两者的差异，
并讨论差异是否集中在特定商品或特定收入组。
\end{remark}

\subsection{第 6 步：福利计算的三个必答问题}

在写出任何 CV/EV 数字之前，有三个问题必须先有答案：

\begin{enumerate}
\item \textbf{对象是谁？} 是个体消费者、某类家庭的代表，还是"代表性消费者"？
      三者的福利含义不同，且只有第一种有明确的偏好基础。
\item \textbf{价格变化的路径是什么？} 单一价格变动时 CV/EV 无路径依赖（命题~\ref{prop:no-path}），
      多价格同时变动时也无——但如果用 $\Delta\mathrm{CS}$ 做近似，路径必须交代。
\item \textbf{分配效应有多大？} 若政策同时改变相对价格与收入分配，
      代表性消费者的 CV 可能与人头加总的福利变化符号相反（\S 6.7 的注）。
      至少应报告分收入组的 CV，看方向是否一致。
\end{enumerate}

\subsection{一个可复用的检验清单}

把本章的技术条件整理成一张可以直接贴进论文附录的清单：

\begin{center}
\begin{tabular}{lll}
\toprule
检验对象 & 原假设 & 拒绝的含义\\
\midrule
零次齐次 & $\sum_j\gamma_{ij}=0$ & 存在货币幻觉或价格指数误设\\
Walras 律 & $\sum_ip_ix_i=w$ & 储蓄、耐用品存量或数据口径问题\\
对称性 & $S=S^{\top}$（或 $\gamma_{ij}=\gamma_{ji}$） & 异质、动态或测量误差\\
负半定性 & $S$ 的特征值 $\le0$ & 补偿需求向上倾斜，理论被实质违反\\
GARP & 无循环严格偏好 & 数据不可由任何局部非饱和效用函数生成\\
\bottomrule
\end{tabular}
\end{center}
\figexplan{五个检验的对象与解读。前两行是会计恒等式，拒绝通常指示数据问题而非理论问题；后三行才是实质检验。注意检验的次序有讲究：应先查前两行（数据是否自洽），再查对称与曲率（理论是否成立），最后查 GARP（是否需要放弃参数化）。顺序颠倒会导致把数据错误误判为理论失败。}

\section{本章小结}
% ============================================================
% 第六讲章末小结：逻辑链、一页检查清单与常见陷阱
% ============================================================

本章把前五章的单向推导倒过来走了一遍：从可观测的需求出发，经可积性条件恢复支出函数，
经显示偏好处理有限数据，最后回到货币度量。把十个小节串成一条逻辑链：

\begin{enumerate}
\item \textbf{问题}（\S 6.1）：反问题有三层——函数层面的可积性、数据层面的可理性化、福利层面的货币度量。
      层次不同，所需信息与结论强度不同。
\item \textbf{信息含量}（\S 6.1）：从偏好到需求是多对一的，因此最多恢复到序数结构；
      可识别的是曲率（$S$），不可识别的是水平（$u$ 的标度）。
\item \textbf{会计约束}（\S 6.2）：齐次性与 Walras 律给出 Engel 加总、Cournot 加总，
      并自动蕴含 $Sp=0$、$p^{\top}S=0$、$\operatorname{rank}S\le n-1$。
      这些是恒等式，不是理论检验的对象。
\item \textbf{对称性的内容}（\S 6.3）：$n=2$ 时对称性自动成立（命题~\ref{prop:n2-auto}），
      实质内容退化为单条曲率条件 $\alpha\le0$；$n\ge3$ 时对称性才成为可拒绝的约束（命题~\ref{prop:n3-substantive}）。
      负半定与对称是两条独立条件。
\item \textbf{可积性}（\S 6.4）：Shephard 引理把问题写成 Pfaff 系统 $\partial e/\partial p_i=x_i(p,e)$；
      其 Frobenius 条件恰为 $S$ 对称（命题~\ref{prop:frobenius}）；
      加上 $S$ 负半定即得 Hurwicz--Uzawa 定理（定理~\ref{thm:integrability}），
      与必要性合起来构成"可由偏好生成 $\iff$ 对称 $+$ 负半定"的充要刻画。
\item \textbf{显式恢复}（\S 6.5）：CD、CES、LES 的 $S$ 都可以显式写出；
      LES 与 CD 同构（$S^{\mathrm{LES}}=A\cdot S^{\mathrm{CD}}(a,1)$）；
      任何显式效用函数推出的系统都自动可积，理论检验的机会随之消失。
\item \textbf{非参数}（\S 6.6）：有限数据下 GARP 是可理性化的充要条件（Afriat）；
      SARP $\Rightarrow$ GARP，但 GARP 与 WARP 互不包含；
      两商品时 WARP、SARP、GARP 三者等价（Rose, 1958）。
\item \textbf{加总}（\S 6.7）：总量需求一般依赖财富分配，不构成代表性消费者；
      加总 Slutsky 矩阵等于个体之和加一个分配项（命题~\ref{prop:S-agg}）；
      Gorman 极化形式 $e_i=c_i(p)+d(p)u_i$ 是精确加总的充分条件，代价是效用必须可作基数比较。
\item \textbf{福利}（\S 6.8）：$\mathrm{CV}=e(p^1,u^0)-w$、$\mathrm{EV}=w-e(p^0,u^1)$；
      单一价格变动时 $\mathrm{EV}\le|\Delta\mathrm{CS}|\le\mathrm{CV}$（命题~\ref{prop:sandwich}）；
      多价格变动时 CV/EV 无路径依赖而 $\Delta\mathrm{CS}$ 有（命题~\ref{prop:no-path}）。
\item \textbf{执行}（\S 6.9）：六步流水线；三种拒绝（异质、动态、测量误差）无法由数据区分；
      写 CV/EV 之前必须回答对象、路径、分配效应三个问题。
\end{enumerate}

\noindent 从需求数据到福利的完整逻辑链，按顺序自检：

\begin{enumerate}
\item \textbf{信息边界}：需求只给边际信息，效用只到序数；能恢复的是支出函数，不是效用函数。
\item \textbf{可观测构造}：$S_{ij}=\partial x_i/\partial p_j+x_j\partial x_i/\partial w$，全部由可观测导数组成。
\item \textbf{自动性质}：$Sp=0$、$p^{\top}S=0$、$\operatorname{rank}S\le n-1$——由齐次与 Walras 律直接给出，不能当检验。
\item \textbf{维数效应}：$n=2$ 时 $S$ 必对称，实质约束只剩曲率；$n\ge3$ 时对称性才有内容。
\item \textbf{可积条件}：Pfaff 系统的 Frobenius 条件 $\iff S$ 对称；$S$ 负半定 $\iff e$ 对 $p$ 凹。
\item \textbf{充分性}：$S$ 对称 $+$ 负半定 $\Rightarrow$ 存在 $e$ 与偏好（Hurwicz--Uzawa）；与必要性合成充要刻画。
\item \textbf{序数残差}：$e$ 只在选定效用标度后唯一；换标度不改变 $S$、需求与福利变化的货币值。
\item \textbf{参数系统}：显式效用函数推出的 $S$ 自动可积；CD 份额恒定、CES 弹性恒定、LES 分两阶段。
\item \textbf{非参数}：GARP $\iff$ 存在局部非饱和效用函数；$e^*=1$ 当且仅当 GARP；两商品时三公理等价。
\item \textbf{加总}：$S^{\mathrm{agg}}=\sum_iS^i+$ 分配项；Gorman 极化形式保证精确加总。
\item \textbf{福利}：$\mathrm{CV}$ 与 $\mathrm{EV}$ 由 $e$ 定义，故无路径依赖；$\Delta\mathrm{CS}$ 只在单一价格变动时良定义。
\item \textbf{解释边界}：拒绝可能来自异质、动态或测量误差；CV/EV 不可跨人相加；代表性与分配效应须交代。
\end{enumerate}

\medskip
\noindent\textbf{三个常见陷阱。}
其一，把"行和为零"当作 Slutsky 理论的验证——它是齐次性与 Walras 律的推论，
在任何估计中都会被自动满足，拒绝它只说明数据或代码有问题，不说明理论有问题。
其二，在两商品情形下做对称性检验并报告"通过"——对称性在那里是恒等式（命题~\ref{prop:n2-auto}），
检验的功效为零；真正的检验必须至少三种商品。
其三，用总量数据算 CV 并直接称为"社会福利变化"——这句话同时跨过了加总、
效用可比与代表性三层未验证的门槛；规范的写法是把结论限定在"代表性消费者的补偿变差"，
并报告分组的稳健性。

\end{document}
